\documentclass[11pt]{article}
\usepackage[margin=1in]{geometry}
\usepackage[T1]{fontenc}
\usepackage[utf8]{inputenc}
\usepackage[french]{babel}
\usepackage{amsmath,amssymb,amsthm}
\usepackage[colorlinks=true,linkcolor=blue,urlcolor=blue]{hyperref}
\newtheorem{problem}{Problème}
\newenvironment{refs}{\par\small\noindent\emph{Références :}\begin{itemize}\setlength\itemsep{0pt}}{\end{itemize}\normalsize}
\title{Hors de la Caverne \\ \Large Équations aux dérivées partielles}
\author{}
\date{}
\begin{document}
\maketitle
\noindent\emph{Des problèmes, pas de réponses.} Chaque problème ci-dessous est posé
sans solution. Les références indiquent où trouver les connaissances nécessaires.
\medskip


\section*{Licence}

\begin{problem}[{Les ondes progressives résolvent l'équation des ondes --- Licence}]
\textbf{PDE-01.} \textbf{Problème.} Soit $c > 0$ et soient $f, g : \mathbb{R} \to \mathbb{R}$
deux fois dérivables.
\begin{enumerate}
\item[(a)] Montrer que $u(x,t) = f(x-ct) + g(x+ct)$ satisfait l'équation
des ondes $u_{tt} = c^2 u_{xx}$ : toute onde se propageant vers la gauche ou vers la droite à
la vitesse $c$, quelle qu'en soit la forme, est une solution.
\item[(b)] Vérifier que $u(x,t) = e^{-kt}\sin x$ satisfait
l'équation de la chaleur $u_t = k u_{xx}$.
\item[(c)] Déterminer tous les couples de constantes $(a,b)$ pour
lesquels $e^{at}\sin(bx)$ résout l'équation de la chaleur, et tous ceux pour lesquels il résout
l'équation des ondes.
\end{enumerate}

\end{problem}

\noindent C'est la porte d'entrée : rien de plus que la règle de dérivation composée à
plusieurs variables n'est requis, et pourtant la partie (a) est essentiellement la découverte de
d'Alembert en 1747, à savoir que la solution générale de l'équation de la corde est une
superposition de deux ondes progressives. Vérifier qu'une fonction donnée résout une EDP est la
première compétence du métier --- la seconde, bien plus difficile, est d'aller dans l'autre
sens.

\begin{refs}
\item Contexte : Équation des ondes --- Wikipédia (\url{https://fr.wikipedia.org/wiki/%C3%89quation_des_ondes})
\item Contexte : Équation de la chaleur --- Wikipédia (\url{https://fr.wikipedia.org/wiki/%C3%89quation_de_la_chaleur})
\item Lecture : W. A. Strauss, \textit{Partial Differential Equations: An Introduction}, chap. 1--2
\end{refs}

\bigskip

\begin{problem}[{Premières rencontres avec les fonctions harmoniques --- Licence}]
\textbf{PDE-02.} \textbf{Problème.} Une fonction $u$ à dérivées secondes continues est
\textit{harmonique} si $\Delta u = u_{xx} + u_{yy} = 0$.
\begin{enumerate}
\item[(a)] Vérifier que
$u(x,y) = \log(x^2+y^2)$ est harmonique sur le plan privé de l'origine.
\item[(b)] Vérifier
que $x^2 - y^2$ et $x^3 - 3xy^2$ sont harmoniques, et démontrer que pour tout
$n \ge 1$ la partie réelle et la partie imaginaire de $(x+iy)^n$ sont harmoniques.
\item[(c)] Calculer la valeur moyenne de $x^2 - y^2$ sur le cercle unité centré à l'origine et la
comparer à la valeur au centre. Formuler une conjecture.
\end{enumerate}

\end{problem}

\noindent L'équation de Laplace $\Delta u = 0$ régit les potentiels gravitationnels et
électrostatiques dans le vide, la chaleur en régime stationnaire et les écoulements
incompressibles irrotationnels ; Euler et Laplace l'ont rencontrée dans les années 1780 et elle est
depuis lors le carrefour de l'analyse. La partie (b) est le premier aperçu du mariage entre
fonctions harmoniques et analyse complexe, et la partie (c) annonce la propriété de la moyenne
(PDE-08), le fait le plus utile qui soit à propos des fonctions harmoniques.

\begin{refs}
\item Contexte : Fonction harmonique --- Wikipédia (\url{https://fr.wikipedia.org/wiki/Fonction_harmonique})
\item Contexte : Équation de Laplace --- Wikipédia (\url{https://fr.wikipedia.org/wiki/%C3%89quation_de_Laplace})
\item Lecture : W. A. Strauss, \textit{Partial Differential Equations: An Introduction}, chap. 6
\end{refs}

\bigskip

\begin{problem}[{L'équation de transport et ses caractéristiques --- Licence}]
\textbf{PDE-03.} \textbf{Problème.}
\begin{enumerate}
\item[(a)] Soit $c$ une constante. Montrer que toute solution
$C^1$ de $u_t + c u_x = 0$ est constante le long de chaque droite
$x - ct = \text{const}$, et en déduire que la solution vérifiant $u(x,0) = \varphi(x)$
est $u(x,t) = \varphi(x - ct)$.
\item[(b)] Résoudre l'équation à coefficient variable
$u_t + x u_x = 0$ avec $u(x,0) = \varphi(x)$ : trouver les courbes le long desquelles $u$
est constante, et écrire explicitement la solution.
\item[(c)] Démontrer la vitesse de propagation finie
pour (a) : si $\varphi$ est nulle hors de $[-M, M]$, alors $u(\cdot,t)$ est nulle
hors de $[-M - |c|t,\ M + |c|t]$.
\end{enumerate}

\end{problem}

\noindent L'équation de transport (advection) est la plus simple EDP qui ne soit pas une
équation différentielle ordinaire déguisée, et elle contient déjà l'idée centrale des équations
hyperboliques : l'information voyage le long de courbes particulières appelées
\textit{caractéristiques}, et le long d'elles l'EDP se réduit à une équation différentielle
ordinaire. La méthode remonte à Euler, Monge et Cauchy ; poussée plus loin, elle résout les
équations du premier ordre pleinement non linéaires (PDE-12) et explique pourquoi les solutions de
l'équation de Burgers se brisent (PDE-13).

\begin{refs}
\item Contexte : Méthode des caractéristiques --- Wikipédia (\url{https://fr.wikipedia.org/wiki/M%C3%A9thode_des_caract%C3%A9ristiques})
\item Contexte : Advection --- Wikipédia (\url{https://fr.wikipedia.org/wiki/Advection})
\item Lecture : W. A. Strauss, \textit{Partial Differential Equations: An Introduction}, chap. 1--2
\item Cours : MIT OCW 18.303 Linear Partial Differential Equations: Analysis and Numerics (\url{https://ocw.mit.edu/courses/18-303-linear-partial-differential-equations-analysis-and-numerics-fall-2014/})
\end{refs}

\bigskip

\begin{problem}[{Établir l'équation de la chaleur --- Licence}]
\textbf{PDE-04.} \textbf{Problème.} Une barre mince et homogène a pour température $u(x,t)$,
pour chaleur massique $c$, pour masse volumique $\rho$ et pour conductivité thermique
$K_0$, toutes constantes. On prend deux postulats physiques : l'énergie thermique d'un segment
$[a,b]$ vaut $\int_a^b c\rho u\,dx$ (conservation de l'énergie), et la chaleur s'écoule du
chaud vers le froid à un taux proportionnel au gradient de température, avec un flux
$q = -K_0 u_x$ (loi de Fourier).
\begin{enumerate}
\item[(a)] En écrivant le taux de variation de l'énergie dans un segment \textit{arbitraire}
$[a,b]$ comme le flux à travers ses extrémités, et en utilisant l'arbitraire de $a$
et $b$, établir $u_t = k u_{xx}$ avec $k = K_0/(c\rho)$. Justifier soigneusement le
passage de \og{} l'intégrale sur tout $[a,b]$ est nulle \fg{} à \og{} la fonction intégrée est
nulle \fg{}.
\item[(b)] Recommencer en trois dimensions avec le théorème de la divergence pour obtenir
$u_t = k\Delta u$.
\item[(c)] Modifier la dérivation pour inclure une source de chaleur interne
$Q(x,t)$ par unité de volume.
\end{enumerate}

\end{problem}

\noindent Fourier a mené essentiellement cette dérivation dans son mémoire de 1807 à
l'Académie de Paris, développé ensuite dans la \textit{Théorie analytique de la chaleur}
(1822). Le schéma --- une loi de conservation plus une loi de comportement (flux proportionnel
à un gradient) donne une équation de diffusion --- se retrouve partout : la loi de Fick donne la
diffusion chimique, la loi d'Ohm donne le courant électrique, et le même argument en génétique des
populations et en finance donne les équations de Fisher et de Black--Scholes.
Apprendre à établir une EDP, et pas seulement à en résoudre une, c'est la moitié du sujet.

\begin{refs}
\item Contexte : Équation de la chaleur --- Wikipédia (\url{https://fr.wikipedia.org/wiki/%C3%89quation_de_la_chaleur})
\item Contexte : Conduction thermique --- Wikipédia (\url{https://fr.wikipedia.org/wiki/Conduction_thermique})
\item Lecture : R. Haberman, \textit{Applied Partial Differential Equations with Fourier Series and Boundary Value Problems}, chap. 1
\end{refs}

\bigskip

\begin{problem}[{Séparation des variables sur un intervalle --- Licence}]
\textbf{PDE-05.} \textbf{Problème.} On considère $u_t = k u_{xx}$ sur $0 < x < L$
avec $u(0,t) = u(L,t) = 0$ et $u(x,0) = \varphi(x)$.
\begin{enumerate}
\item[(a)] Trouver toutes les solutions
de la forme $u = X(x)T(t)$ satisfaisant les conditions au bord, obtenant les modes
$\sin(n\pi x/L)\,e^{-k(n\pi/L)^2 t}$.
\item[(b)] Les superposer et identifier les
coefficients $b_n$ comme les coefficients de Fourier en sinus de $\varphi$.
\item[(c)] Supposer
$\varphi$ bornée et intégrable. Démontrer que pour tout $t > 0$ la série
converge, peut être dérivée terme à terme, et résout l'équation --- la solution est donc
infiniment régulière pour $t > 0$ même si $\varphi$ est irrégulière.
\item[(d)] Montrer que la solution décroît vers $0$ comme $e^{-k(\pi/L)^2 t}$ : le mode le plus lent
l'emporte.
\end{enumerate}

\end{problem}

\noindent C'est la méthode de Fourier lui-même, de 1807. Lagrange et le comité de
l'Académie ont regimbé devant l'affirmation qu'une température initiale arbitraire puisse se
développer en sinus --- un doute de haute lignée (voir PDE-06) --- et le prix qu'ils ont fini par
décerner à Fourier en 1812 s'accompagnait d'une mention notant un manque de rigueur. Rendre (c)
précis est exactement le genre de travail qui a transformé le doute en théorèmes, et la valeur
propre qui apparaît en (d) revient en héroïne de PDE-17.

\begin{refs}
\item Contexte : Séparation des variables --- Wikipédia (\url{https://fr.wikipedia.org/wiki/S%C3%A9paration_des_variables})
\item Lecture : R. Haberman, \textit{Applied Partial Differential Equations with Fourier Series and Boundary Value Problems}, chap. 2--3
\item Cours : MIT OCW 18.303 Linear Partial Differential Equations: Analysis and Numerics (\url{https://ocw.mit.edu/courses/18-303-linear-partial-differential-equations-analysis-and-numerics-fall-2014/})
\end{refs}

\bigskip

\begin{problem}[{La corde vibrante et la formule de d'Alembert --- Licence}]
\textbf{PDE-06.} \textbf{Problème.} (a) En factorisant l'opérateur des ondes en
$(\partial_t - c\,\partial_x)(\partial_t + c\,\partial_x)$, établir la formule de d'Alembert
pour $u_{tt} = c^2 u_{xx}$ sur toute la droite avec $u(x,0) = \varphi(x)$
et $u_t(x,0) = \psi(x)$ :
\[ u(x,t) = \tfrac{1}{2}\big[\varphi(x-ct) + \varphi(x+ct)\big]
+ \frac{1}{2c}\int_{x-ct}^{x+ct} \psi(s)\,ds. \]
(b) Résoudre le cas de la corde pincée : une corde fixée en $x = 0$ et $x = L$, lâchée sans
vitesse initiale dans une forme triangulaire, en prolongeant les données à la droite entière en
fonctions impaires $2L$-périodiques. Décrire géométriquement le mouvement. (c) Le $\varphi$
triangulaire n'est pas deux fois dérivable ; en quel sens (b) \og{} résout \fg{}-elle donc l'équation des
ondes ? Proposer un sens précis et le défendre.
\end{problem}

\noindent D'Alembert a publié l'équation des ondes et cette solution en 1747 --- et a
aussitôt restreint les formes initiales à celles données par une seule expression analytique.
Euler a objecté (1748) : une corde pincée a un coin, et la physique se moque de savoir si la forme
possède une formule. Daniel Bernoulli (1753) a affirmé au contraire que tout mouvement est une
superposition de modes sinusoïdaux, ce que d'Alembert comme Euler jugeaient absurde pour des
formes générales. Tous trois avaient partiellement raison, et la bataille sur le sens des mots
\og{} fonction \fg{} et \og{} solution \fg{} a couvé soixante ans, jusqu'à ce que Fourier prenne le parti de
Bernoulli avec une méthode (PDE-05), que Dirichlet fasse de la convergence un théorème, et que le
concept moderne de fonction émerge. La partie (c) est la porte des solutions faibles
(PDE-14) --- la réponse du XXe siècle à l'objection d'Euler.

\begin{refs}
\item Contexte : Équation de d'Alembert (et sa formule) --- Wikipédia (\url{https://fr.wikipedia.org/wiki/%C3%89quation_de_d%27Alembert})
\item Contexte : Onde sur une corde vibrante --- Wikipédia (\url{https://fr.wikipedia.org/wiki/Onde_sur_une_corde_vibrante})
\item Lecture : W. A. Strauss, \textit{Partial Differential Equations: An Introduction}, chap. 2 et 4
\item Cours : MIT OCW 18.152 Introduction to Partial Differential Equations (\url{https://ocw.mit.edu/courses/18-152-introduction-to-partial-differential-equations-fall-2011/})
\end{refs}

\bigskip

\begin{problem}[{Séries de Fourier et phénomène de Gibbs --- Licence}]
\textbf{PDE-07.} \textbf{Problème.} Soit $f$ le signal carré : $f(x) = 1$ pour
$0 < x < \pi$ et $f(x) = -1$ pour $-\pi < x < 0$.
\begin{enumerate}
\item[(a)] Calculer sa
série de Fourier
\[ \frac{4}{\pi}\Big(\sin x + \frac{\sin 3x}{3} + \frac{\sin 5x}{5} + \cdots\Big) \]
et déterminer, avec démonstration, sa limite ponctuelle en tout $x$, points de saut
compris.
\item[(b)] Soit $S_N$ la somme partielle jusqu'à la fréquence $N$. Montrer que
$S_N$ a un maximum local près de $x = \pi/N$ et démontrer que
\[ S_N(\pi/N) \;\longrightarrow\; \frac{2}{\pi}\int_0^{\pi} \frac{\sin t}{t}\,dt
\qquad (N \to \infty). \]
Évaluer ce nombre numériquement et en conclure que les sommes partielles dépassent le saut
d'environ 9 \% de sa hauteur --- une erreur qui se rapproche du saut mais ne diminue jamais.
\item[(c)] En déduire que la convergence n'est uniforme sur aucun intervalle contenant
$0$, et expliquer pourquoi cela ne contredit pas (a).
\end{enumerate}

\end{problem}

\noindent Le dépassement a été trouvé par Henry Wilbraham en 1848 puis oublié ; il a
refait surface en 1898 lorsque l'analyseur harmonique d'Albert Michelson --- une machine qui
sommait 80 modes de Fourier --- n'a cessé de dessiner de petites cornes aux sauts, et que Michelson
a écrit à \textit{Nature} en soupçonnant un défaut. J. Willard Gibbs a expliqué en 1899 que ces
cornes sont mathématiques, non mécaniques. Le phénomène est une leçon permanente sur la différence
entre convergence simple et convergence uniforme, et il compte encore : c'est l'artefact de
sonnerie aux bords francs en traitement du signal et de l'image.

\begin{refs}
\item Contexte : Phénomène de Gibbs --- Wikipédia (\url{https://fr.wikipedia.org/wiki/Ph%C3%A9nom%C3%A8ne_de_Gibbs})
\item Contexte : Convergence des séries de Fourier --- Wikipédia (en anglais) (\url{https://en.wikipedia.org/wiki/Convergence_of_Fourier_series})
\item Lecture : W. A. Strauss, \textit{Partial Differential Equations: An Introduction}, chap. 5
\item Lecture : T. W. Körner, \textit{Fourier Analysis}, Cambridge University Press
\end{refs}

\bigskip

\begin{problem}[{La propriété de la moyenne --- Licence}]
\textbf{PDE-08.} \textbf{Problème.} Soit $u$ une fonction harmonique sur un ouvert
$\Omega \subset \mathbb{R}^2$.
\begin{enumerate}
\item[(a)] Démontrer la propriété de la moyenne : pour tout disque
fermé de $\Omega$ de centre $p$ et de rayon $r$, la moyenne de $u$ sur le
cercle frontière vaut $u(p)$.
\item[(b)] En déduire que la moyenne sur le disque plein
vaut également $u(p)$.
\item[(c)] Démontrer la réciproque : si $u$ est continue sur $\Omega$
et vérifie la propriété de la moyenne sur les cercles en tout point et pour tout petit rayon,
alors $u$ est harmonique.
\item[(d)] Utiliser (a) pour démontrer le théorème de Liouville dans le plan : une
fonction harmonique sur tout $\mathbb{R}^2$ et bornée supérieurement et inférieurement est
constante.
\end{enumerate}

\end{problem}

\noindent Gauss connaissait la propriété de la moyenne dans les années 1840, dans ses
travaux sur la théorie du potentiel. C'est la définition de l'équité selon l'analyste : une
fonction harmonique prend en chaque point exactement la moyenne des valeurs de ses voisines, ce
qui explique pourquoi les fonctions harmoniques décrivent des équilibres --- et pourquoi le gain
espéré d'un marcheur aléatoire est harmonique, un pont entre EDP et probabilités. Presque tout ce
qu'il y a de profond sur les fonctions harmoniques (principes du maximum, régularité, inégalités
de Harnack) découle de cette seule propriété.

\begin{refs}
\item Contexte : Fonction harmonique --- Wikipédia (\url{https://fr.wikipedia.org/wiki/Fonction_harmonique})
\item Lecture : L. C. Evans, \textit{Partial Differential Equations}, chap. 2
\item Cours : MIT OCW 18.152 Introduction to Partial Differential Equations (\url{https://ocw.mit.edu/courses/18-152-introduction-to-partial-differential-equations-fall-2011/})
\end{refs}

\bigskip

\begin{problem}[{La brachistochrone --- Licence}]
\textbf{PDE-09.} \textbf{Problème.} Jean Bernoulli, \textit{Acta Eruditorum},
juin 1696, s'adressant \og{} aux mathématiciens les plus perspicaces du monde entier \fg{} :
\textit{étant donné deux points A et B dans un plan vertical, trouver la courbe le long de laquelle
une particule, partant du repos en A et glissant sous l'effet de la pesanteur sans frottement,
atteint B en un temps minimal.} (a) En utilisant la conservation de l'énergie (vitesse
$\sqrt{2gy}$ après une chute de hauteur $y$), exprimer le temps de parcours comme la
fonctionnelle
\[ T[y] = \int_0^{x_1} \sqrt{\frac{1 + y'(x)^2}{2g\,y(x)}}\;dx. \]
(b) Établir l'équation d'Euler--Lagrange pour une fonctionnelle $\int F(y,y')\,dx$, et
montrer que lorsque $F$ ne fait pas intervenir $x$ explicitement elle admet l'intégrale
première $F - y'\,\partial F/\partial y' = \text{const}$ (identité de Beltrami).
(c) En conclure que la courbe minimisante vérifie $y\,(1 + y'^2) = \text{const}$,
et vérifier que la cycloïde --- le chemin $x = R(\theta - \sin\theta)$,
$y = R(1 - \cos\theta)$ tracé par un point d'un cercle qui roule --- satisfait cette
équation. (d) Montrer que la réponse n'est ni la droite ni un arc de cercle.
\end{problem}

\noindent C'est le problème qui a créé le calcul des variations.
Des solutions sont venues de Leibniz, de l'Hôpital, de Jacques Bernoulli, de Jean lui-même --- et une
anonyme de Newton, composée dit-on en une nuit après une journée à la Monnaie ;
Jean en a reconnu l'auteur \og{} comme le lion à sa griffe \fg{}. Galilée avait deviné un
arc de cercle en 1638 et se trompait. Euler (1744) et Lagrange (1755) ont transformé
l'astuce en méthode générale, et l'idée que la nature minimise une intégrale ---
action, énergie, temps --- est devenue l'un des principes organisateurs de la physique et de
la théorie des EDP elle-même (voir PDE-16).

\begin{refs}
\item Contexte : Courbe brachistochrone --- Wikipédia (\url{https://fr.wikipedia.org/wiki/Courbe_brachistochrone})
\item Contexte : Calcul des variations --- Wikipédia (\url{https://fr.wikipedia.org/wiki/Calcul_des_variations})
\item Lecture : I. M. Gelfand et S. V. Fomin, \textit{Calculus of Variations}, chap. 1
\end{refs}

\bigskip

\begin{problem}[{Les trois espèces d'équation du second ordre --- Licence, short}]
\textbf{PDE-21.} \textbf{Problème.} Pour l'équation à coefficients constants
$a\,u_{xx} + 2b\,u_{xy} + c\,u_{yy} + (\text{termes du premier ordre}) = 0$, démontrer qu'un
changement de variables linéaire réel réduit la partie principale à exactement l'une des formes
$u_{\xi\xi} + u_{\eta\eta}$, $u_{\xi\xi} - u_{\eta\eta}$ ou $u_{\xi\xi}$,
selon le signe du discriminant $b^2 - ac$, et que ce signe est
inchangé par toute substitution de ce type. En conclure que l'équation de Laplace, l'équation
des ondes et l'équation de la chaleur sont les trois formes normales --- elliptique, hyperbolique,
parabolique --- et qu'aucun changement de coordonnées ne transforme l'une en l'autre.
\end{problem}

\noindent Le vocabulaire est emprunté aux coniques, où le même
discriminant sépare ellipse, hyperbole et parabole ; le transfert aux
équations différentielles date du XIXe siècle et a été rendu décisif par
Hadamard, qui a montré que le type dicte quelles conditions annexes sont légitimes.
Une équation elliptique veut des données tout autour d'un bord, une équation hyperbolique veut
position et vitesse initiales, et une équation parabolique ne s'écoule que dans un seul sens du
temps. Demander à une équation hyperbolique de résoudre un problème de Dirichlet n'est pas
seulement difficile --- c'est mal posé.

\begin{refs}
\item Contexte : Équation aux dérivées partielles elliptique --- Wikipédia (\url{https://fr.wikipedia.org/wiki/%C3%89quation_aux_d%C3%A9riv%C3%A9es_partielles_elliptique})
\item Contexte : Équation aux dérivées partielles hyperbolique --- Wikipédia (\url{https://fr.wikipedia.org/wiki/%C3%89quation_aux_d%C3%A9riv%C3%A9es_partielles_hyperbolique})
\item Lecture : W. A. Strauss, \textit{Partial Differential Equations: An Introduction}, chap. 1
\end{refs}

\bigskip

\begin{problem}[{L'unicité par l'énergie --- Licence, short}]
\textbf{PDE-22.} \textbf{Problème.} Pour l'équation des ondes $u_{tt} = c^2 u_{xx}$
sur $0 \le x \le L$ avec $u(0,t) = u(L,t) = 0$, démontrer que
\[ E(t) = \frac12\int_0^L \bigl(u_t^2 + c^2 u_x^2\bigr)\,dx \]
est indépendante de $t$, et en déduire que le problème mixte a au plus une solution
$C^2$. Démontrer le même énoncé d'unicité pour les conditions de Neumann
$u_x(0,t) = u_x(L,t) = 0$, et expliquer où les conditions au bord
sont utilisées.
\end{problem}

\noindent C'est la méthode de l'énergie sous sa forme la plus pure : dériver une
quantité quadratique, intégrer par parties, voir les termes de bord s'annuler, et obtenir
l'unicité sans jamais écrire une solution. À comparer avec l'unicité qui
découle de la formule de d'Alembert (PDE-06) --- la formule n'est disponible que dans des
géométries particulières, tandis que l'argument d'énergie survit aux coefficients variables,
aux dimensions supérieures et aux problèmes non linéaires, où il devient le premier
outil standard du sujet.

\begin{refs}
\item Contexte : Équation des ondes --- Wikipédia (\url{https://fr.wikipedia.org/wiki/%C3%89quation_des_ondes})
\item Lecture : W. A. Strauss, \textit{Partial Differential Equations: An Introduction}, chap. 2
\item Lecture : L. C. Evans, \textit{Partial Differential Equations}, chap. 2
\end{refs}

\bigskip

\begin{problem}[{Le principe de Duhamel --- Licence}]
\textbf{PDE-23.} \textbf{Problème.} Le principe de Duhamel affirme qu'une équation d'évolution linéaire
forcée, à données nulles, se résout en superposant les évolutions libres du
forçage, traité comme un flux d'impulsions instantanées. Rendre cela précis dans trois
cadres.
\begin{enumerate}
\item[(a)] Démontrer le prototype pour les équations différentielles ordinaires : si $A$ est une
matrice constante, l'unique solution de
$u'(t) - Au(t) = f(t)$, $u(0) = 0$, est
\[ u(t) = \int_0^t e^{(t-s)A} f(s)\,ds. \]
Expliquer pourquoi la fonction intégrée est \og{} l'évolution libre, pendant le temps restant
$t-s$, du forçage délivré à l'instant $s$ \fg{}.
\item[(b)] Soit $\Phi(x,t)$ le noyau de la chaleur de PDE-11. Démontrer que
\[ u(x,t) = \int_0^t \int_{\mathbb{R}} \Phi(x-y,\,t-s)\,f(y,s)\,dy\,ds \]
résout $u_t - k u_{xx} = f$ sur $\mathbb{R}\times(0,\infty)$ avec $u(x,0)=0$,
pour $f$ lisse à support compact. Attention : la dérivation naïve
sous le signe intégral échoue en $s = t$, et réparer cette étape est le cœur de la
démonstration.
\item[(c)] Démontrer la version pour l'équation des ondes : la solution de $u_{tt} - c^2 u_{xx} = f$
avec $u(x,0) = u_t(x,0) = 0$ est
\[ u(x,t) = \frac{1}{2c}\int_0^t \int_{x-c(t-s)}^{x+c(t-s)} f(y,s)\,dy\,ds, \]
en appliquant la formule de d'Alembert aux problèmes auxiliaires à données impulsionnelles.
\item[(d)] Utiliser (c) avec la séparation des variables pour étudier une corde de longueur $L$
excitée par $f(x,t) = \sin(n\pi x/L)\sin(\omega t)$. Montrer que la réponse reste
bornée lorsque $\omega \ne n\pi c/L$ mais croît linéairement en $t$ lorsque la fréquence
d'excitation coïncide avec une fréquence propre --- la résonance, établie et non affirmée.
\end{enumerate}

\end{problem}

\noindent Jean-Marie Duhamel a introduit l'idée vers 1830 en étudiant
la conduction de la chaleur dans un corps dont la température de surface varie dans le temps : il
a remplacé la condition au bord variant continûment par une somme de petits sauts brusques, résolu
chacun, et additionné. Le geste est si naturel qu'il ressemble aujourd'hui à une définition, mais
il encode un vrai théorème --- que l'opérateur solution d'une équation linéaire autonome forme un
semi-groupe, et que le problème forcé en est la convolution avec le forçage. Au
XXe siècle la même formule est devenue la \textit{solution mild} des équations
d'évolution abstraites, et c'est le point de départ standard pour démontrer qu'une EDP non
linéaire est localement bien posée : mettre la non-linéarité dans le rôle de $f$, et
chercher un point fixe.

\begin{refs}
\item Contexte : Principe de Duhamel --- Wikipédia (en anglais) (\url{https://en.wikipedia.org/wiki/Duhamel%27s_principle})
\item Lecture : W. A. Strauss, \textit{Partial Differential Equations: An Introduction}, chap. 3
\item Lecture : L. C. Evans, \textit{Partial Differential Equations}, chap. 2
\end{refs}

\bigskip

\begin{problem}[{Le tambour rectangulaire --- Licence}]
\textbf{PDE-24.} \textbf{Problème.} Une membrane occupe le rectangle
$R = [0,a]\times[0,b]$, est encastrée le long de son bord et vibre selon
$u_{tt} = c^2\,\Delta u$ avec $u = 0$ sur $\partial R$.
\begin{enumerate}
\item[(a)] Par séparation des variables, montrer que le problème aux valeurs propres
$-\Delta \varphi = \lambda\varphi$ avec $\varphi = 0$ sur $\partial R$ a pour
solutions $\varphi_{mn}(x,y) = \sin\frac{m\pi x}{a}\sin\frac{n\pi y}{b}$ pour des
entiers $m,n \ge 1$, avec
\[ \lambda_{mn} = \pi^2\Bigl(\frac{m^2}{a^2} + \frac{n^2}{b^2}\Bigr), \]
de sorte que les fréquences du tambour sont $c\sqrt{\lambda_{mn}}$.
\item[(b)] Démontrer que les $\varphi_{mn}$ sont deux à deux orthogonales dans $L^2(R)$, et
écrire la solution de données initiales $u(\cdot,0) = g$, $u_t(\cdot,0) = h$ comme une
série double de Fourier en sinus, en donnant explicitement les coefficients.
\item[(c)] Montrer que, contrairement à la corde de PDE-06, les fréquences ne sont en général pas des
multiples entiers de la plus basse : exhiber deux d'entre elles dont le rapport est irrationnel.
C'est le sens précis dans lequel un tambour rend un son sans hauteur définie.
\item[(d)] Prendre le carré, $a = b$. Montrer que $\lambda_{mn}$ est dégénérée --- l'espace propre
est de dimension supérieure à un --- exactement lorsque $m^2 + n^2$ admet plus
d'une représentation comme somme ordonnée de deux carrés positifs, et en donner un
exemple. Décrire comment l'ensemble nodal de $\alpha\varphi_{mn} + \beta\varphi_{nm}$
change quand $(\alpha,\beta)$ parcourt un espace propre dégénéré.
\end{enumerate}

\end{problem}

\noindent La partie (d) est l'endroit où un fait d'arithmétique --- quels entiers sont
sommes de deux carrés, et de combien de façons --- devient audible et visible.
Ernst Chladni avait déjà rendu les ensembles nodaux visibles en 1787 en frottant à l'archet des
plaques métalliques saupoudrées de sable ; ses figures ont tant impressionné Napoléon que
l'Académie de Paris a mis un prix au concours pour les expliquer, remporté après trois tentatives
par Sophie Germain en 1816. Le dénombrement des valeurs propres dans cet exemple,
$\#\{\lambda_{mn} \le \Lambda\} \sim \frac{|R|}{4\pi}\Lambda$, est le premier cas de la loi de
Weyl, et il dit que le tambour connaît sa propre aire --- ce qui est exactement le fil repris dans
PDE-17.

\begin{refs}
\item Contexte : Équation de Helmholtz --- Wikipédia (\url{https://fr.wikipedia.org/wiki/%C3%89quation_de_Helmholtz})
\item Contexte : Loi de Weyl --- Wikipédia (en anglais) (\url{https://en.wikipedia.org/wiki/Weyl_law})
\item Contexte : Ernst Chladni (et ses figures) --- Wikipédia (\url{https://fr.wikipedia.org/wiki/Ernst_Chladni})
\item Lecture : R. Haberman, \textit{Applied Partial Differential Equations}, chap. 7
\end{refs}

\bigskip

\begin{problem}[{Le tambour circulaire et les fonctions de Bessel --- Licence}]
\textbf{PDE-31.} \textbf{Problème.} Une membrane couvre le disque de rayon $a$, est encastrée le long de son
bord, et obéit à $u_{tt} = c^2 \Delta u$ avec $u = 0$ sur la frontière. En coordonnées
polaires
\[ \Delta u = u_{rr} + \frac{1}{r}\,u_r + \frac{1}{r^2}\,u_{\theta\theta}. \]
\begin{enumerate}
\item[(a)] Séparer les variables, $u = R(r)\,\Theta(\theta)\,T(t)$. Montrer que l'uniformité en
$\theta$ impose $\Theta(\theta) = \cos n\theta$ ou $\sin n\theta$ avec $n$ entier
positif ou nul, et que $R$ vérifie
\[ r^2 R'' + r R' + (k^2 r^2 - n^2)\,R = 0, \]
qui devient l'équation de Bessel d'ordre $n$ dans la variable $s = kr$. Expliquer quel
principe permet d'écarter la seconde fonction de Bessel et de garder $R(r) = J_n(kr)$.
\item[(b)] Imposer le bord encastré : montrer que $k a$ doit être un zéro positif $j_{n,m}$ de
$J_n$, de sorte que les fréquences sont $c\,j_{n,m}/a$. Démontrer que les fonctions propres
obtenues sont deux à deux orthogonales dans le $L^2$ du disque muni de l'élément d'aire
$r\,dr\,d\theta$ --- mettre d'abord l'équation radiale sous la forme auto-adjointe
$\bigl(r R'\bigr)' + \bigl(k^2 r - n^2/r\bigr) R = 0$ --- et écrire la solution de déplacement et
de vitesse initiaux donnés comme une série double.
\item[(c)] Démontrer que le tambour est sans hauteur définie. À partir de $J_0' = -J_1$ et
$(x J_1)' = x J_0$, montrer que les zéros positifs de $J_0$ et $J_1$ s'entrelacent
strictement. Puis, avec $j_{0,1} \approx 2{,}405$ et $j_{1,1} \approx 3{,}832$, montrer que la
deuxième fréquence n'est pas un multiple entier de la première --- contrairement à la corde de
PDE-06, dont les harmoniques sont des multiples exacts du fondamental --- et expliquer ce que cela
signifie pour l'oreille.
\item[(d)] Décrire l'ensemble nodal du mode $(n,m)$ : montrer qu'il est constitué de $n$ diamètres
et de $m-1$ cercles concentriques, et identifier ces courbes aux lignes le long desquelles le
sable s'accumule sur une plaque circulaire frottée à l'archet. Régler ensuite les dégénérescences :
montrer que toute valeur propre avec $n \ge 1$ est de multiplicité au moins deux, et démontrer
que l'hypothèse de Bourget --- que $J_n$ et $J_{n+m}$ n'ont aucun zéro positif commun, établie
par Carl Ludwig Siegel en 1929 --- implique qu'il n'y en a pas d'autres. Comparer au tambour carré
de PDE-24, où les dégénérescences abondent.
\end{enumerate}

\end{problem}

\noindent Euler a traité la membrane circulaire vibrante en coordonnées cylindriques
vers 1780 et a introduit pour cela les séries entières que nous appelons aujourd'hui $J_n$ ; le
mémoire de Bessel de 1824 sur les perturbations planétaires a donné à ces fonctions leur première
étude systématique et leur nom, et la \textit{Theory of Sound} de Rayleigh (1877) en a fait l'outil
de travail de l'acoustique. Ernst Chladni avait déjà rendu la partie (d) visible en 1787 en
promenant un archet le long du bord d'une plaque couverte de sable. Le contraste avec le rectangle
est tout l'objet du problème : la même équation, un bord changé, et l'arithmétique du spectre est
entièrement différente --- ce qui est exactement pourquoi le spectre peut encoder la géométrie
(PDE-17).

\begin{refs}
\item Contexte : Vibration d'une membrane circulaire --- Wikipédia (en anglais) (\url{https://en.wikipedia.org/wiki/Vibration_of_a_circular_membrane})
\item Contexte : Fonction de Bessel --- Wikipédia (\url{https://fr.wikipedia.org/wiki/Fonction_de_Bessel})
\item Lecture : W. A. Strauss, \textit{Partial Differential Equations: An Introduction}, chap. 10
\item Lecture : R. Haberman, \textit{Applied Partial Differential Equations}, chap. 7
\end{refs}

\bigskip

\begin{problem}[{Séparer l'équation de Laplace sur la boule --- Licence, short}]
\textbf{PDE-32.} \textbf{Problème.} Écrire l'équation de Laplace en coordonnées sphériques et
la séparer sous la forme $u(r,\theta,\phi) = R(r)\,Y(\theta,\phi)$ ; démontrer que la constante
de séparation doit valoir $n(n+1)$ pour un entier naturel $n$ si $Y$ doit être lisse sur toute
la sphère --- ramener l'équation en $\theta$ à l'équation de Legendre dans la variable
$\cos\theta$ --- que le facteur radial est alors une combinaison de $r^{n}$ et $r^{-n-1}$, et
donc que la solution du problème de Dirichlet sur la boule unité avec pour donnée au bord une
harmonique sphérique $Y_n$ est $u = r^{n} Y_n(\theta,\phi)$. En déduire que $r^n Y_n$ est un
polynôme harmonique homogène de degré $n$, et que pour de telles données la valeur au centre est
égale à la moyenne des valeurs au bord --- la propriété de la moyenne de PDE-08, retrouvée en
dimension trois.
\end{problem}

\noindent Legendre (1782) et Laplace (1785) ont inventé ces fonctions pour calculer
l'attraction gravitationnelle d'une planète quasi sphérique, et pendant un siècle on les a
appelées coefficients de Laplace ; le \textit{Treatise on Natural Philosophy} de Kelvin et Tait en a
fait l'équipement standard des physiciens. Elles restent les coordonnées dans lesquelles le monde
rond se rapporte : les coefficients multipolaires du fond diffus cosmologique, le modèle harmonique
du champ de gravité terrestre et la partie angulaire de toute orbitale atomique sont des
harmoniques sphériques. La séparation menée ici est aussi la source de l'équation de Legendre
elle-même --- la sphère est son lieu d'origine, non un exercice arbitraire sur les séries.

\begin{refs}
\item Contexte : Harmonique sphérique --- Wikipédia (\url{https://fr.wikipedia.org/wiki/Harmonique_sph%C3%A9rique})
\item Contexte : Équation de Laplace --- Wikipédia (\url{https://fr.wikipedia.org/wiki/%C3%89quation_de_Laplace})
\item Lecture : W. A. Strauss, \textit{Partial Differential Equations: An Introduction}, chap. 10
\end{refs}

\bigskip

\begin{problem}[{Théorie de Sturm--Liouville : pourquoi la séparation des variables fonctionne --- Licence}]
\textbf{PDE-33.} \textbf{Problème.} Un \textit{problème de Sturm--Liouville régulier} sur $[a,b]$ est
\[ -\bigl(p(x)\,u'\bigr)' + q(x)\,u = \lambda\,w(x)\,u, \]
avec $p \in C^1$, $q, w$ continues, $p > 0$ et $w > 0$ sur $[a,b]$, assorti de
conditions au bord séparées $\alpha_1 u(a) + \alpha_2 u'(a) = 0$ et
$\beta_1 u(b) + \beta_2 u'(b) = 0$, où $(\alpha_1,\alpha_2)$ et $(\beta_1,\beta_2)$
sont non nuls.
\begin{enumerate}
\item[(a)] Démontrer l'identité de Lagrange : avec $Lu = (p u')' - q u$, l'expression $v\,Lu - u\,Lv$
est la dérivée de $p\,(u'v - u v')$. Intégrer pour obtenir la formule de Green, et démontrer que
les termes de bord s'annulent pour deux fonctions quelconques satisfaisant les conditions séparées
--- l'opérateur est symétrique.
\item[(b)] En déduire les trois faits fondamentaux : toute valeur propre est réelle ; les fonctions
propres associées à des valeurs propres distinctes sont orthogonales pour le poids $w$ ; et
chaque valeur propre est simple. (Pour ce dernier point, utiliser l'unicité pour le problème de
Cauchy en $x = a$.)
\item[(c)] Établir le quotient de Rayleigh : si $u$ est une fonction propre de valeur propre
$\lambda$, exprimer $\lambda$ comme le rapport de $\int_a^b \bigl(p (u')^2 + q u^2\bigr)\,dx$
plus des termes de bord à $\int_a^b w u^2\,dx$. En déduire une minoration du spectre lorsque
$q \ge 0$ et que les conditions au bord sont $u(a) = u(b) = 0$, et démontrer la
caractérisation variationnelle de la plus petite valeur propre comme minimum du quotient.
\item[(d)] Reconnaître votre propre travail passé. Identifier $p, q, w$ et l'intervalle pour : la corde
vibrante de PDE-05 ; l'équation radiale du tambour circulaire de PDE-31, où
$p(r) = w(r) = r$ s'annule en $r = 0$ ; et l'équation de Legendre sur $[-1,1]$, où
$p = 1-x^2$ s'annule aux deux extrémités. Énoncer le théorème de complétude pour le problème
régulier --- les fonctions propres forment une base orthogonale de l'espace $L^2$ à poids --- et
expliquer ce qui change dans les deux cas \textit{singuliers}, en particulier pourquoi une exigence
de bornitude y remplace une condition au bord.
\end{enumerate}

\end{problem}

\noindent Charles-François Sturm et Joseph Liouville, amis et codirecteurs du
\textit{Journal de Mathématiques Pures et Appliquées}, ont publié en 1836--1837 une série d'articles
liés qui ont fondé la théorie spectrale des opérateurs différentiels : une infinité de valeurs
propres réelles s'en allant vers l'infini, des théorèmes d'oscillation et de comparaison pour les
fonctions propres, et un théorème de développement --- soixante-dix ans avant que Hilbert n'en
fournisse le cadre général. C'est le théorème qui fait de la séparation des variables autre chose
qu'un tour heureux : les modes qu'elle produit sont des fonctions propres d'un opérateur
symétrique, donc elles sont orthogonales, les coefficients se calculent par intégration, et le
développement d'un état initial arbitraire est garanti convergent. Tout ce qu'on a reproché à
Fourier d'avoir supposé en 1807 (PDE-05) est démontré ici.

\begin{refs}
\item Contexte : Théorie de Sturm-Liouville --- Wikipédia (\url{https://fr.wikipedia.org/wiki/Th%C3%A9orie_de_Sturm-Liouville})
\item Contexte : Quotient de Rayleigh --- Wikipédia (\url{https://fr.wikipedia.org/wiki/Quotient_de_Rayleigh})
\item Lecture : R. Haberman, \textit{Applied Partial Differential Equations}, chap. 5
\item Lecture : R. Courant \& D. Hilbert, \textit{Methods of Mathematical Physics}, vol. 1
\end{refs}

\bigskip


\section*{Master}

\begin{problem}[{Principes du maximum et unicité --- Master}]
\textbf{PDE-10.} \textbf{Problème.} (a) Démontrer le principe du maximum faible : une
fonction harmonique sur un domaine borné $\Omega$ et continue sur son adhérence
atteint son maximum sur la frontière $\partial\Omega$. (Considérer
$u + \varepsilon|x|^2$ pour $\varepsilon > 0$ petit.) (b) Démontrer la version
parabolique : une solution de l'équation de la chaleur sur un rectangle espace-temps atteint son
maximum sur la frontière parabolique (bas et côtés). (c) Utiliser (a) pour démontrer que le
problème de Dirichlet $\Delta u = f$ dans $\Omega$, $u = g$ sur $\partial\Omega$
a au plus une solution, et utiliser (b) pour démontrer l'unicité et la dépendance continue
en les données pour le problème mixte de l'équation de la chaleur. (d) Démontrer le
principe du maximum fort pour les fonctions harmoniques via la propriété de la moyenne : si
$u$ atteint son maximum en un point intérieur d'un domaine connexe, alors $u$ est
constante.
\end{problem}

\noindent Les principes du maximum sont la bête de somme de l'elliptique et du
parabolique : des démonstrations d'unicité, de stabilité et de comparaison qui n'exigent aucune
formule explicite de solution. Physiquement, ils disent qu'un corps sans source interne de chaleur
ne peut développer un point chaud. La technique va de ces énoncés classiques jusqu'au lemme de
Hopf (1927) et à l'analyse géométrique du XXIe siècle --- les arguments de flot de Ricci
de Hamilton et de Perelman sont, au fond, des principes du maximum sophistiqués.

\begin{refs}
\item Contexte : Principe du maximum --- Wikipédia (\url{https://fr.wikipedia.org/wiki/Principe_du_maximum_(%C3%A9quations_aux_d%C3%A9riv%C3%A9es_partielles)})
\item Lecture : L. C. Evans, \textit{Partial Differential Equations}, chap. 2
\item Lecture : M. H. Protter et H. F. Weinberger, \textit{Maximum Principles in Differential Equations}
\end{refs}

\bigskip

\begin{problem}[{Le noyau de la chaleur --- Master}]
\textbf{PDE-11.} \textbf{Problème.} (a) Montrer que si $u(x,t)$ résout
$u_t = k u_{xx}$, alors $u(\lambda x, \lambda^2 t)$ aussi, pour tout $\lambda$.
Chercher une solution de la forme $t^{-1/2}F(x/\sqrt{t})$ d'intégrale totale $1$,
établir l'équation différentielle ordinaire vérifiée par $F$, et la résoudre pour obtenir le
noyau de la chaleur
\[ \Phi(x,t) = \frac{1}{\sqrt{4\pi k t}}\; e^{-x^2/(4kt)}. \]
(b) Démontrer que pour $\varphi$ continue bornée, la convolution
$u(x,t) = \int_{\mathbb{R}} \Phi(x-y,t)\,\varphi(y)\,dy$ résout l'équation de la chaleur
pour $t > 0$ et $u(x,t) \to \varphi(x_0)$ quand $(x,t) \to (x_0, 0^+)$.
(c) Montrer que si $\varphi \ge 0$ n'est pas identiquement nulle et s'annule hors d'un
ensemble borné, alors $u(x,t) > 0$ pour tout $x$ et tout $t > 0$ : la chaleur
se propage à vitesse infinie. (d) Montrer que $u(\cdot,t)$ est infiniment dérivable
pour tout $t > 0$, et énoncer le noyau en dimension $n$. Comparer (c) et (d)
à la vitesse finie et à la conservation des singularités de l'équation des ondes.
\end{problem}

\noindent Le noyau de la chaleur est la \textit{solution élémentaire} : la réponse
à une unité de chaleur concentrée en un point, à partir de laquelle toute solution se construit
par superposition --- l'idée introduite par George Green pour l'électrostatique en 1828. La
même gaussienne est la densité de transition du mouvement brownien (Einstein 1905, Wiener
1923), ce qui fait du noyau de la chaleur la charnière entre EDP et probabilités ; en géométrie
moderne, son comportement en temps court encode la courbure et même la topologie (la
démonstration du théorème de l'indice d'Atiyah--Singer par le noyau de la chaleur).

\begin{refs}
\item Contexte : Noyau de la chaleur --- Wikipédia (\url{https://fr.wikipedia.org/wiki/Noyau_de_la_chaleur})
\item Contexte : Solution élémentaire --- Wikipédia (en anglais) (\url{https://en.wikipedia.org/wiki/Fundamental_solution})
\item Lecture : L. C. Evans, \textit{Partial Differential Equations}, chap. 2
\end{refs}

\bigskip

\begin{problem}[{Caractéristiques pour les équations non linéaires du premier ordre --- Master}]
\textbf{PDE-12.} \textbf{Problème.} On considère une équation du premier ordre pleinement
non linéaire $F(\nabla u, u, x) = 0$ sur un domaine de $\mathbb{R}^n$.
\begin{enumerate}
\item[(a)] Établir le
système caractéristique : en écrivant $p = \nabla u$ le long d'une courbe $x(s)$, montrer que
si $u$ est une solution lisse on peut choisir $x(s)$ de sorte que $x$, $z = u(x(s))$
et $p$ satisfassent le système fermé d'équations différentielles ordinaires
\[ \dot{x} = \nabla_p F, \qquad \dot{z} = p \cdot \nabla_p F, \qquad
\dot{p} = -\nabla_x F - p\,\frac{\partial F}{\partial z}, \]
de sorte que l'EDP se résout en intégrant des équations différentielles ordinaires à partir des
données au bord.
\item[(b)] Vérifier que pour les
équations quasi linéaires le système se réduit aux caractéristiques de PDE-03 et
PDE-13.
\item[(c)] L'appliquer à l'équation eikonale $|\nabla u| = 1$ avec $u = 0$ sur
la frontière d'un domaine borné régulier : montrer que les caractéristiques sont des rayons
rectilignes normaux à la frontière et que, près de la frontière, $u(x)$ est la distance à
celle-ci.
\item[(d)] Montrer par un exemple (un domaine non convexe, ou une composante de bord
intérieure) que les rayons finissent par se croiser, et qu'aucune solution $C^1$ ne peut exister
au-delà du croisement --- la régularité meurt là où les caractéristiques se percutent.
\end{enumerate}

\end{problem}

\noindent Les équations différentielles caractéristiques des équations non linéaires ont
été trouvées par Lagrange, Monge et Charpit dans les années 1770--1780 et perfectionnées par Cauchy.
Hamilton a refondu l'optique géométrique de cette manière dans les années 1830 : les rayons
lumineux sont les caractéristiques de l'équation eikonale, et les fronts d'onde ses lignes de
niveau --- et la même structure, appliquée à la mécanique, est devenue l'équation de
Hamilton--Jacobi. Le croisement des rayons en (d) est une caustique, la courbe brillante dans votre
tasse de café, et l'échec des solutions classiques en ce lieu est ce que les solutions de
viscosité (Crandall--Lions, 1983) ont été inventées pour réparer.

\begin{refs}
\item Contexte : Méthode des caractéristiques --- Wikipédia (\url{https://fr.wikipedia.org/wiki/M%C3%A9thode_des_caract%C3%A9ristiques})
\item Contexte : Équation eikonale --- Wikipédia (\url{https://fr.wikipedia.org/wiki/%C3%89quation_eikonale})
\item Lecture : L. C. Evans, \textit{Partial Differential Equations}, chap. 3
\end{refs}

\bigskip

\begin{problem}[{L'équation de Burgers : les solutions régulières se brisent --- Master}]
\textbf{PDE-13.} \textbf{Problème.} On considère $u_t + u u_x = 0$ avec une donnée
initiale régulière bornée $u(x,0) = \varphi(x)$.
\begin{enumerate}
\item[(a)] Montrer que les caractéristiques sont des
droites de pente déterminée par $\varphi$, et que toute solution $C^1$
satisfait l'équation implicite $u = \varphi(x - ut)$.
\item[(b)] Démontrer : si
$\varphi' \ge 0$ partout, une solution $C^1$ globale existe.
\item[(c)] Démontrer : si
$\varphi'(x_0) < 0$ pour un certain $x_0$, la solution $C^1$ existe exactement jusqu'au
temps $T^* = 1/\sup(-\varphi')$, et $\sup|u_x| \to \infty$ quand $t \to T^*$.
(Suivre $w = u_x$ le long d'une caractéristique ; elle vérifie $\dot{w} = -w^2$.)
(d) Calculer $T^*$ pour $\varphi(x) = e^{-x^2}$, et décrire géométriquement ce que
fait le graphe de $u(\cdot,t)$ quand $t \to T^*$.
\end{enumerate}

\end{problem}

\noindent C'est la catastrophe du gradient : une onde dont les parties hautes voyagent
plus vite que les parties basses se redresse jusqu'à ce que la pente soit verticale --- après quoi
il n'existe plus aucune solution classique. Riemann a compris le mécanisme en dynamique des gaz
en 1860 ; l'équation porte le nom de Jan Burgers, qui l'a employée dans les années 1940 comme
modèle de turbulence. La rupture n'est pas un défaut du modèle --- les gaz réels forment des ondes
de choc --- c'est un défaut de la notion classique de solution, et la réparer est l'affaire de
PDE-14.

\begin{refs}
\item Contexte : Équation de Burgers --- Wikipédia (\url{https://fr.wikipedia.org/wiki/%C3%89quation_de_Burgers})
\item Lecture : L. C. Evans, \textit{Partial Differential Equations}, chap. 3
\item Lecture : G. B. Whitham, \textit{Linear and Nonlinear Waves}, chap. 2
\end{refs}

\bigskip

\begin{problem}[{Solutions faibles, chocs et entropie --- Master}]
\textbf{PDE-14.} \textbf{Problème.} Écrire l'équation de Burgers sous forme conservative
$u_t + (u^2/2)_x = 0$ et appeler $u$ \textit{solution faible} si
\[ \int_0^\infty\!\!\int_{\mathbb{R}} \Big( u\,\phi_t + \frac{u^2}{2}\,\phi_x \Big)\,dx\,dt
+ \int_{\mathbb{R}} u(x,0)\,\phi(x,0)\,dx = 0 \]
pour toute fonction test $\phi$ lisse à support compact --- une définition qui ne réclame
aucune dérivée de $u$.
\begin{enumerate}
\item[(a)] Montrer que les solutions classiques sont des solutions
faibles, et qu'une solution faible régulière par morceaux présentant un saut le long d'une courbe
$x = \xi(t)$ doit satisfaire la relation de Rankine--Hugoniot : la vitesse du choc
$\xi'$ est égale au saut de $u^2/2$ divisé par le saut de $u$ --- ici, la
moyenne $(u_- + u_+)/2$ des deux côtés.
\item[(b)] Résoudre le problème de Riemann avec
$u = 0$ à gauche de l'origine et $u = 1$ à droite, deux fois : une fois avec un
choc voyageant à la vitesse $1/2$, une fois avec un éventail de détente $u = x/t$ remplissant
le secteur. Vérifier que \textit{les deux} sont des solutions faibles des mêmes données ---
l'unicité est perdue.
\item[(c)] Énoncer la condition d'entropie de Lax (les caractéristiques doivent entrer dans un
choc, non en sortir) et déterminer laquelle des solutions de (b) elle admet et laquelle
elle interdit.
\item[(d)] Montrer que pour les données inversées ($1$ à gauche, $0$ à
droite) le choc satisfait bien la condition d'entropie. Expliquer le nom : laquelle de
ces solutions survivrait à l'ajout d'un petit terme de viscosité ?
\end{enumerate}

\end{problem}

\noindent Les solutions faibles répondent à la crise de PDE-13 et, enfin, au camp d'Euler
dans la querelle de la corde de 1747 : une \og{} solution \fg{} n'a pas besoin d'être dérivable, seulement
d'être correcte face à tout observateur lisse. La relation de saut est due à Rankine (1870)
et Hugoniot (1887), travaillant sur la dynamique des gaz ; la non-unicité de (b) fut le
scandale que les conditions d'entropie (Lax 1957, Oleïnik, Kroujkov) ont résolu, restaurant
une unique solution physiquement pertinente. Ce cercle d'idées --- intégrer contre des
fonctions test, puis se battre pour retrouver l'unicité --- est le modèle de bien des EDP
modernes, des distributions au problème de Navier--Stokes (PDE-18).

\begin{refs}
\item Contexte : Formulation faible (solution faible) --- Wikipédia (\url{https://fr.wikipedia.org/wiki/Formulation_faible})
\item Contexte : Relations de Rankine-Hugoniot --- Wikipédia (\url{https://fr.wikipedia.org/wiki/Relations_de_Rankine-Hugoniot})
\item Lecture : L. C. Evans, \textit{Partial Differential Equations}, chap. 3 et 11
\item Article : P. D. Lax, \og{} Hyperbolic systems of conservation laws II \fg{} (1957), Communications on Pure and Applied Mathematics 10
\end{refs}

\bigskip

\begin{problem}[{Une porte d'entrée vers les espaces de Sobolev --- Master}]
\textbf{PDE-15.} \textbf{Problème.} On dit que $v$ est la \textit{dérivée faible} de
$u$ sur $(0,1)$ si $\int u\varphi' = -\int v\varphi$ pour toute $\varphi$ lisse à
support compact, et l'on note $H^1(0,1)$ l'espace des $u$ de carré intégrable
à dérivée faible de carré intégrable, muni de la norme
$\|u\|^2 = \int u^2 + \int (u')^2$.
\begin{enumerate}
\item[(a)] Montrer que la dérivée faible est unique (lorsqu'elle
existe) et la calculer pour $u(x) = |x - \tfrac12|$.
\item[(b)] Démontrer l'injection
de Sobolev en dimension un : toute $u \in H^1(0,1)$ admet un représentant
continu, et $|u(x) - u(y)| \le \|u'\|_{L^2}\,|x-y|^{1/2}$.
\item[(c)] Démontrer
l'inégalité de Poincaré pour $u \in H^1(0,1)$ avec $u(0) = u(1) = 0$ :
$\|u\|_{L^2} \le C\,\|u'\|_{L^2}$, et déterminer la meilleure constante $C$.
(La comparer à la première valeur propre de PDE-05(d).) (d) Montrer que l'injection dans les
fonctions continues échoue en dimension deux : vérifier que
$u(x) = \log\log(1 + 1/|x|)$ sur le disque unité est de norme $H^1$ finie mais
non bornée.
\end{enumerate}

\end{problem}

\noindent Les espaces de Sobolev (S. L. Sobolev, années 1930) sont l'habitat naturel des
solutions faibles : des espaces complets où \og{} l'énergie \fg{} est la norme, de sorte que les arguments
de compacité et de passage à la limite qui échouent dans les espaces de fonctions régulières
réussissent. Que la meilleure constante de (c) soit une valeur propre n'a rien d'un hasard --- c'est
la caractérisation variationnelle du spectre, et toute la théorie moderne de l'existence pour les
équations elliptiques (Lax--Milgram, Galerkine, méthodes directes) tourne sur des inégalités comme
(b) et (c). La partie (d) marque la frontière où la dimension commence à compter, le seuil des
théorèmes d'injection de Sobolev.

\begin{refs}
\item Contexte : Espace de Sobolev --- Wikipédia (\url{https://fr.wikipedia.org/wiki/Espace_de_Sobolev})
\item Lecture : L. C. Evans, \textit{Partial Differential Equations}, chap. 5
\item Lecture : H. Brezis, \textit{Functional Analysis, Sobolev Spaces and Partial Differential Equations}
\end{refs}

\bigskip

\begin{problem}[{Le principe de Dirichlet et l'objection de Weierstrass --- Master}]
\textbf{PDE-16.} \textbf{Problème.} Le principe de Dirichlet affirme que la
solution de $\Delta u = 0$ dans $\Omega$ avec $u = g$ sur $\partial\Omega$ est
le minimiseur de l'énergie de Dirichlet $E[w] = \int_\Omega |\nabla w|^2$ parmi toutes les
$w$ égales à $g$ sur la frontière. (a) Démontrer un sens : si une fonction $C^2$
$u$ minimise $E$ parmi les concurrentes $C^2$ de mêmes valeurs au bord,
alors $\Delta u = 0$ dans $\Omega$. (Calculer la variation première
$\frac{d}{d\varepsilon} E[u + \varepsilon\varphi]$.) (b) Démontrer la réciproque : une
fonction harmonique ayant les valeurs au bord données minimise $E$. (c) Voici maintenant
l'objection de Weierstrass --- le minimum pourrait ne pas exister. Étudier
\[ J[y] = \int_{-1}^{1} x^2\,y'(x)^2\,dx \]
sur les fonctions $y$ de classe $C^1$ par morceaux avec $y(-1) = -1$ et $y(1) = 1$.
Démontrer que la borne inférieure de $J$ vaut $0$ (construire une suite minimisante
explicite) et qu'aucune fonction admissible ne l'atteint. (d) Expliquer précisément quelle étape
du raisonnement \og{} $E$ est minorée, donc elle a un minimiseur \fg{} est fautive, et pourquoi (a)
seule ne démontre donc rien sur l'existence pour le problème de Dirichlet.
\end{problem}

\noindent Green, Gauss, Thomson et surtout Riemann --- qui a donné au principe le nom de
son maître Dirichlet et a bâti sur lui son analyse complexe --- tenaient tous l'existence d'un
minimiseur pour évidente : l'énergie est positive, donc quelque chose doit bien atteindre sa plus
petite valeur. Le contre-exemple de Weierstrass en 1870 (la partie (c) est le sien) a montré que
l'argument est un sophisme, et pendant trente ans le théorème de représentation conforme de
Riemann est resté en suspens. Hilbert a réhabilité le principe vers 1900--1904 en démontrant
honnêtement l'existence --- la \og{} méthode directe \fg{} --- un sauvetage qui a exigé d'inventer les bons
espaces de fonctions et les bons arguments de compacité, c'est-à-dire une grande part de l'analyse
moderne. La morale, qu'une borne inférieure n'est pas un minimum tant qu'on n'a pas démontré
qu'elle est atteinte, est l'une des leçons les plus tranchantes des mathématiques.

\begin{refs}
\item Contexte : Principe de Dirichlet --- Wikipédia (\url{https://fr.wikipedia.org/wiki/Principe_de_Dirichlet})
\item Contexte : Méthode directe du calcul des variations --- Wikipédia (en anglais) (\url{https://en.wikipedia.org/wiki/Direct_method_in_the_calculus_of_variations})
\item Lecture : R. Courant, \textit{Dirichlet's Principle, Conformal Mapping, and Minimal Surfaces}
\item Lecture : L. C. Evans, \textit{Partial Differential Equations}, chap. 8
\end{refs}

\bigskip

\begin{problem}[{Peut-on entendre la forme d'un tambour ? --- Master}]
\textbf{PDE-17.} \textbf{Problème.} Les valeurs propres de Dirichlet d'un domaine plan
$\Omega$ sont les nombres $\lambda$ pour lesquels $\Delta u + \lambda u = 0$ a une
solution non nulle s'annulant sur $\partial\Omega$ --- physiquement, les carrés des
fréquences d'une peau de tambour de la forme de $\Omega$.
\begin{enumerate}
\item[(a)] Calculer les valeurs propres et
les fonctions propres du rectangle $[0,a] \times [0,b]$.
\item[(b)] Démontrer la loi de Weyl pour
les rectangles : le nombre $N(\lambda)$ de valeurs propres $\le \lambda$ vérifie
$N(\lambda)/\lambda \to ab/(4\pi)$ quand $\lambda \to \infty$. (Compter les points du
réseau dans un quart d'ellipse.) En conclure qu'un auditeur qui entend toutes les
fréquences d'un tambour rectangulaire peut en retrouver l'aire.
\item[(c)] L'inégalité de
Faber--Krahn (1923--1925, démontrant une conjecture de Rayleigh) affirme que parmi les domaines
d'aire donnée, le disque a la plus petite première valeur propre, et qu'il est le seul
minimiseur. En l'admettant, démontrer qu'un domaine plan ayant exactement le même spectre de
Dirichlet qu'un certain disque doit être ce disque --- on \textit{peut} entendre la forme d'un
tambour circulaire.
\item[(d)] La question de Kac en 1966 demande si deux domaines plans non congruents peuvent partager
tout leur spectre. Énoncer précisément ce que Gordon, Webb et Wolpert ont démontré en
1992, et le concilier avec (b) et (c).
\end{enumerate}

\end{problem}

\noindent L'article de Mark Kac \og{} Can one hear the shape of a drum? \fg{} (1966,
dont le titre reprend un mot d'esprit de Lipman Bers) a rendu célèbre la géométrie spectrale.
Milnor avait déjà trouvé des tores isospectraux non isométriques en dimension 16 (1964), mais le
plan a résisté jusqu'à ce que Carolyn Gordon, David Webb et Scott Wolpert produisent deux
polygones non congruents --- en forme d'hélice, constructibles en papier --- de spectres
identiques : en général, on ne peut pas entendre la forme d'un tambour. Pourtant le spectre
détermine bien l'aire (Weyl 1911), le périmètre (Pleijel), et le fait que le tambour soit un
disque ; savoir si l'on peut entendre la forme d'un tambour \textit{lisse et convexe} est encore
ouvert.

\begin{refs}
\item Contexte : Entendre la forme d'un tambour --- Wikipédia (en anglais) (\url{https://en.wikipedia.org/wiki/Hearing_the_shape_of_a_drum})
\item Contexte : Loi de Weyl --- Wikipédia (en anglais) (\url{https://en.wikipedia.org/wiki/Weyl_law})
\item Article : M. Kac, \og{} Can One Hear the Shape of a Drum? \fg{} (1966), American Mathematical Monthly 73
\item Article : C. Gordon, D. Webb, S. Wolpert, \og{} One cannot hear the shape of a drum \fg{} (1992), Bulletin of the American Mathematical Society 27
\end{refs}

\bigskip

\begin{problem}[{Le noyau de Poisson sur le disque --- Master, short}]
\textbf{PDE-25.} \textbf{Problème.} Démontrer que pour toute fonction $g$ continue sur le
cercle unité, la fonction
\[ u(re^{i\theta}) \;=\; \frac{1}{2\pi}\int_{-\pi}^{\pi}
\frac{1-r^2}{1 - 2r\cos(\theta-\phi) + r^2}\;g(\phi)\,d\phi, \qquad 0 \le r < 1, \]
est harmonique sur le disque unité ouvert et se prolonge continûment au disque fermé avec pour
valeurs au bord $g$, de sorte que le problème de Dirichlet sur le disque est résolu par une
formule explicite. En déduire l'inégalité de Harnack : toute fonction harmonique positive $u$
sur le disque vérifie $\frac{1-r}{1+r}\,u(0) \le u(re^{i\theta}) \le \frac{1+r}{1-r}\,u(0)$.
\end{problem}

\noindent Poisson a écrit le noyau dans les années 1820 ; Schwarz en a fait
la démonstration d'existence moderne. Deux choses en font plus qu'une formule. D'abord, le
noyau est une approximation de l'unité --- il se concentre en $\phi = \theta$ quand
$r \to 1$ --- ce qui explique pourquoi les valeurs au bord sont atteintes, et le même mécanisme
revient partout en analyse. Ensuite, l'inégalité de Harnack est un principe du maximum
quantitatif : elle dit qu'une fonction harmonique positive ne peut être grande quelque part et
minuscule tout à côté, et sous cette forme elle se propage jusqu'à la théorie de la régularité de
De Giorgi--Nash--Moser et aux travaux de Perelman sur le flot de Ricci. Du point de vue
probabiliste, le noyau est la loi de sortie d'un mouvement brownien parti de
$re^{i\theta}$.

\begin{refs}
\item Contexte : Noyau de Poisson --- Wikipédia (\url{https://fr.wikipedia.org/wiki/Noyau_de_Poisson})
\item Contexte : Inégalité de Harnack --- Wikipédia (\url{https://fr.wikipedia.org/wiki/In%C3%A9galit%C3%A9_de_Harnack})
\item Lecture : L. V. Ahlfors, \textit{Complex Analysis}, chap. 4
\item Lecture : L. C. Evans, \textit{Partial Differential Equations}, chap. 2
\end{refs}

\bigskip

\begin{problem}[{La dispersion pour l'équation de Schrödinger libre --- Master, short}]
\textbf{PDE-26.} \textbf{Problème.} Pour l'équation de Schrödinger libre
$i\,u_t = -\Delta u$ sur $\mathbb{R}^n$, démontrer que le flot préserve exactement la
norme $L^2$ tout en obéissant à l'estimation dispersive
\[ \|u(\cdot,t)\|_{L^\infty(\mathbb{R}^n)} \;\le\; \frac{1}{(4\pi|t|)^{n/2}}\;
\|u_0\|_{L^1(\mathbb{R}^n)}, \qquad t \ne 0. \]
Expliquer comment les deux peuvent tenir à la fois --- rien n'est perdu, et pourtant la solution
s'aplatit partout --- et vérifier sur une donnée initiale gaussienne explicite que le taux
$|t|^{-n/2}$ ne peut être amélioré.
\end{problem}

\noindent Les deux énoncés réunis expliquent pourquoi l'équation est dite
dispersive plutôt que dissipative : la masse $L^2$ est conservée mais s'étale,
parce que les ondes planes $e^{i(x\cdot\xi - |\xi|^2 t)}$ de fréquences différentes voyagent
à des vitesses de groupe différentes $2\xi$. La démonstration est un calcul de phase stationnaire
avec le noyau $(4\pi i t)^{-n/2} e^{i|x|^2/4t}$, obtenu à partir du noyau de la chaleur par la
substitution formelle $t \mapsto it$. Les estimations de ce type, affinées par
Strichartz en 1977 en bornes espace-temps, sont ce qui permet de traiter les équations de
Schrödinger non linéaires par des arguments de point fixe, et elles sous-tendent
l'analyse derrière PDE-20.

\begin{refs}
\item Contexte : Équation de Schrödinger --- Wikipédia (\url{https://fr.wikipedia.org/wiki/%C3%89quation_de_Schr%C3%B6dinger})
\item Contexte : Relation de dispersion --- Wikipédia (\url{https://fr.wikipedia.org/wiki/Relation_de_dispersion})
\item Lecture : T. Tao, \textit{Nonlinear Dispersive Equations: Local and Global Analysis}, chap. 2
\end{refs}

\bigskip

\begin{problem}[{Fisher--KPP et la vitesse d'avance d'un gène --- Master}]
\textbf{PDE-27.} \textbf{Problème.} On considère l'équation de Fisher--KPP sur la droite,
\[ u_t = u_{xx} + u(1-u), \qquad 0 \le u \le 1, \]
où $u$ est la fréquence d'un gène avantageux, ou la densité d'une population
envahissante : elle diffuse et croît logistiquement.
\begin{enumerate}
\item[(a)] Chercher un front progressif $u(x,t) = U(x-ct)$ avec $U(-\infty)=1$,
$U(+\infty)=0$. Établir l'équation différentielle $U'' + cU' + U(1-U) = 0$, l'écrire comme
système plan, et classer les points d'équilibre $(1,0)$ et $(0,0)$ ; montrer que
$(0,0)$ est un nœud stable quand $c \ge 2$ et un foyer stable quand $0 < c < 2$.
\item[(b)] En déduire qu'un front avec $U \ge 0$ existe pour tout $c \ge c_* = 2$ et pour
aucun $c$ plus petit, parce qu'une approche en spirale force $U$ à devenir négative --- ce que
la biologie interdit.
\item[(c)] Donner l'argument de détermination linéaire pour le même nombre : linéariser en
$u = 0$ pour obtenir $u_t = u_{xx} + u$, substituer $e^{-\lambda(x-ct)}$, obtenir
$c = \lambda + 1/\lambda$, et minimiser en $\lambda > 0$. Expliquer pourquoi il n'est pas
évident que la vitesse sélectionnée par l'équation non linéaire doive être égale à celle
prédite par sa linéarisation au front d'attaque.
\item[(d)] Énoncer le théorème de Kolmogorov, Petrovski et Piskounov : pour une donnée initiale en
échelon, la solution converge en forme vers le front critique. Puis lire le raffinement de
Bramson, selon lequel la position du front est $m(t) = 2t - \tfrac32\log t + O(1)$, et
expliquer ce que signifie le retard logarithmique.
\end{enumerate}

\end{problem}

\noindent En 1937 R. A. Fisher, se demandant à quelle vitesse un allèle favorable se
répand dans une population étalée le long d'un habitat, a écrit cette équation et calculé
la vitesse $2\sqrt{rD}$ ; la même année, Kolmogorov, Petrovski et Piskounov
publiaient pour elle un théorème de convergence rigoureux --- l'un des premiers travaux sérieux
d'analyse des EDP non linéaires, et un cas frappant de deux traditions très différentes
aboutissant à une même équation. La question de la sélection du front qu'ils ont ouverte est
toujours vivante : quelles non-linéarités sont \og{} tirées \fg{} (vitesse fixée par la linéarisation au
front d'attaque, comme ici) et lesquelles sont \og{} poussées \fg{}. La correction logarithmique de
(d) a été trouvée par Bramson en 1978 en traduisant le problème en mouvement brownien branchant,
et elle reste l'un des ponts les plus cités entre EDP et probabilités.

\begin{refs}
\item Contexte : Équation KPP-Fisher --- Wikipédia (\url{https://fr.wikipedia.org/wiki/%C3%89quation_KPP-Fisher})
\item Contexte : Système de réaction-diffusion --- Wikipédia (\url{https://fr.wikipedia.org/wiki/Syst%C3%A8me_de_r%C3%A9action-diffusion})
\item Article : R. A. Fisher, \og{} The wave of advance of advantageous genes \fg{} (1937), Annals of Eugenics 7 ; A. Kolmogorov, I. Petrovsky, N. Piskunov, \og{} Étude de l'équation de la diffusion\dots{} \fg{} (1937), Bulletin de l'Université d'État à Moscou
\item Lecture : J. D. Murray, \textit{Mathematical Biology I: An Introduction}, chap. 13
\end{refs}

\bigskip

\begin{problem}[{La fonction de Green du laplacien --- Master}]
\textbf{PDE-34.} \textbf{Problème.} Soit $\Phi$ la solution élémentaire du laplacien : une fonction radiale,
harmonique hors de l'origine, normalisée de sorte que $-\Delta \Phi = \delta_0$ --- explicitement
un multiple de $\log|x|$ quand $n = 2$ et de $|x|^{2-n}$ quand $n \ge 3$.
\begin{enumerate}
\item[(a)] Vérifier que $\Phi$ est harmonique sur $\mathbb{R}^n \setminus \{0\}$ et localement
intégrable, et fixer la constante de normalisation. Démontrer ensuite la formule de représentation
de Green pour $u \in C^2(\overline{\Omega})$ sur un domaine borné à frontière régulière :
appliquer la seconde identité de Green sur $\Omega$ privé d'une petite boule autour de $x$ et
faire tendre le rayon vers zéro, obtenant $u(x)$ en fonction de $\Delta u$ dans $\Omega$ et
des valeurs au bord de $u$ et de $\partial u/\partial \nu$.
\item[(b)] La formule n'est pas encore une solution du problème de Dirichlet, car elle utilise des
données au bord dont on ne dispose pas. La réparer : définir la fonction de Green
$G(x,y) = \Phi(y-x) - h^x(y)$, où $h^x$ est harmonique dans $\Omega$ et coïncide avec
$\Phi(\cdot - x)$ sur $\partial\Omega$. Démontrer que $G$ est unique, que
$G(x,y) = G(y,x)$, et que la dérivée normale de $G$ donne une formule de solution explicite
pour $\Delta u = 0$ dans $\Omega$, $u = g$ sur $\partial\Omega$.
\item[(c)] Calculer $G$ dans les deux cas classiques : pour le demi-espace $\{x_n > 0\}$ en
réfléchissant la source à travers la frontière, et pour la boule par l'inversion de Kelvin
$y \mapsto y/|y|^2$ avec un poids adéquat. Retrouver à partir du second calcul le noyau de
Poisson de PDE-25.
\item[(d)] Enfin, résoudre l'équation de Poisson sur tout $\mathbb{R}^n$ : démontrer que
$u = \Phi * f$ satisfait $-\Delta u = f$ pour $f$ deux fois continûment dérivable à support
compact. Expliquer exactement pourquoi dériver deux fois sous le signe intégral est illégitime, et
comment le transfert des dérivées sur $f$ répare l'argument --- la même manœuvre qu'exige la
formule de Duhamel en PDE-23(b). Commenter ce que l'on sait tomber en défaut si $f$ est
seulement continue.
\end{enumerate}

\end{problem}

\noindent George Green, meunier de Nottingham ayant reçu environ une année de scolarité
formelle, a publié son \textit{Essay on the Application of Mathematical Analysis to the Theories of
Electricity and Magnetism} en 1828 par souscription privée --- cinquante et un exemplaires. Il y
introduisait le mot \og{} potentiel \fg{}, les identités qui portent aujourd'hui son nom, et le procédé sur
lequel repose ce problème : résoudre une fois pour une source ponctuelle, puis superposer. L'Essay
est resté ignoré jusqu'à ce que William Thomson, futur lord Kelvin, le découvre en 1845 et le fasse
réimprimer ; Kelvin a aussi fourni la méthode des images et l'inversion utilisée en (c). Le procédé
est l'idée la plus transposable des EDP linéaires --- c'est le noyau de la chaleur de PDE-11, la
réponse impulsionnelle que mesure un ingénieur, et le propagateur de la théorie quantique des
champs.

\begin{refs}
\item Contexte : Fonction de Green --- Wikipédia (\url{https://fr.wikipedia.org/wiki/Fonction_de_Green})
\item Contexte : Méthode des charges images --- Wikipédia (en anglais) (\url{https://en.wikipedia.org/wiki/Method_of_image_charges})
\item Contexte : Transformation de Kelvin --- Wikipédia (en anglais) (\url{https://en.wikipedia.org/wiki/Kelvin_transform})
\item Lecture : L. C. Evans, \textit{Partial Differential Equations}, chap. 2
\end{refs}

\bigskip

\begin{problem}[{Le principe de Huygens et pourquoi les dimensions impaires sont particulières --- Master}]
\textbf{PDE-35.} \textbf{Problème.} Résoudre $u_{tt} = \Delta u$ avec les données $u(\cdot,0) = g$,
$u_t(\cdot,0) = h$, lisses et à support compact, et comparer la réponse selon les
dimensions.
\begin{enumerate}
\item[(a)] En dimension trois, soit $U(x;r,t)$ la moyenne de $u(\cdot,t)$ sur la sphère de
rayon $r$ centrée en $x$. Démontrer l'équation d'Euler--Poisson--Darboux
$U_{tt} = U_{rr} + \frac{2}{r}U_r$, montrer que $rU$ résout l'équation des ondes
unidimensionnelle en $(r,t)$, et en déduire la formule de Kirchhoff exprimant $u(x,t)$ au
moyen des seules moyennes de $g$, $\nabla g$ et $h$ sur la sphère $|y - x| = t$.
\item[(b)] Descendre en dimension deux : considérer une solution dans le plan comme une solution dans
l'espace ne dépendant pas de la troisième coordonnée, et mener la méthode de descente de Hadamard
pour obtenir la formule de Poisson. Observer que l'intégrale obtenue est prise sur tout le disque
plein $|y-x| \le t$, et non sur le cercle qui le borde.
\item[(c)] En conclure le principe de Huygens fort et son échec. En dimension trois, la valeur en
$(x,t)$ ne dépend que des données sur une sphère, si bien qu'une impulsion brève produit un écho
bref puis le silence ; en dimension deux, le disque entier contribue, si bien que toute
perturbation laisse une traîne qui décroît mais ne s'arrête jamais. Énoncer le résultat général ---
le principe fort vaut pour l'équation des ondes en dimension d'espace impaire $n \ge 3$ et
échoue en dimension paire ainsi qu'en dimension un --- et en donner la lecture physique : un étang
continue de résonner, l'air non.
\item[(d)] Problème de compréhension. Expliquer pourquoi la parole et la musique sont possibles dans un
espace de dimension trois et seraient une bouillie brouillée en dimension deux, puis se documenter
sur le problème de Hadamard --- la classification des opérateurs hyperboliques du second ordre qui
satisfont le principe de Huygens fort. Rapporter ce qui est connu, et noter honnêtement que la
classification est généralement décrite comme incomplète.
\end{enumerate}

\end{problem}

\noindent Huygens a bâti l'optique dans son \textit{Traité de la lumière} (1690) à partir
d'ondelettes secondaires émises par chaque point d'un front d'onde --- un argument qui suppose
silencieusement la propagation sans traîne, nettement localisée, que seules les dimensions impaires
fournissent. La formule de Kirchhoff de 1882 en a fait un calcul, et les \textit{Lectures on Cauchy's
Problem} de Hadamard (1923) en ont fait un théorème et ont nommé la méthode de descente. La
morale est répétée par tous les manuels depuis Hadamard : que nous puissions nous entendre les uns
les autres est un fait sur la dimension de l'espace, et le contraste avec la vitesse de propagation
infinie et le lissage instantané de l'équation de la chaleur (PDE-11) est l'illustration la plus
tranchante de cette page de la façon dont le type d'une équation contrôle profondément son
comportement.

\begin{refs}
\item Contexte : Principe de Huygens-Fresnel --- Wikipédia (\url{https://fr.wikipedia.org/wiki/Principe_de_Huygens-Fresnel})
\item Contexte : Méthode de descente de Hadamard --- Wikipédia (en anglais) (\url{https://en.wikipedia.org/wiki/Hadamard%27s_method_of_descent})
\item Lecture : L. C. Evans, \textit{Partial Differential Equations}, chap. 2
\item Lecture : W. A. Strauss, \textit{Partial Differential Equations: An Introduction}, chap. 9
\end{refs}

\bigskip

\begin{problem}[{L'alternative de Fredholm pour un problème aux limites --- Master, short}]
\textbf{PDE-36.} \textbf{Problème.} Pour $f$ continue, démontrer que le problème de Neumann
$-u'' = f$ sur $(0,1)$ avec $u'(0) = u'(1) = 0$ a une solution si et seulement si
$\int_0^1 f\,dx = 0$, et que la solution est alors unique à une constante additive près.
Démontrer l'énoncé analogue pour $-u'' - \pi^2 u = f$ avec $u(0) = u(1) = 0$ --- résoluble
exactement lorsque $\int_0^1 f(x)\sin(\pi x)\,dx = 0$, les solutions étant uniques à des
multiples de $\sin \pi x$ près --- et formuler ce que ces deux exemples illustrent : la
résolubilité vaut précisément lorsque la donnée est orthogonale à toute solution du problème
adjoint homogène, et l'ensemble des solutions est alors une classe latérale de l'espace des
solutions homogènes.
\end{problem}

\noindent Le mémoire d'Ivar Fredholm de 1903 dans \textit{Acta Mathematica} sur les
équations intégrales a produit l'alternative sous sa forme générale --- soit le problème
inhomogène est uniquement résoluble pour toute donnée, soit le problème homogène a un nombre fini
de solutions indépendantes et la résolubilité coûte exactement autant de conditions
d'orthogonalité. C'est le descendant en dimension infinie du théorème du rang ; Hilbert a bâti la
théorie spectrale sur elle, et Riesz et Schauder l'ont refondue en un énoncé sur les opérateurs
compacts. En pratique, c'est la comptabilité de la résonance : la condition d'orthogonalité est ce
qui échoue à une fréquence propre, produisant la croissance linéaire en temps de PDE-23(d), et dans
tout développement perturbatif c'est la condition qui fixe le terme d'ordre suivant.

\begin{refs}
\item Contexte : Alternative de Fredholm --- Wikipédia (\url{https://fr.wikipedia.org/wiki/Alternative_de_Fredholm})
\item Contexte : Condition aux limites de Neumann --- Wikipédia (\url{https://fr.wikipedia.org/wiki/Condition_aux_limites_de_Neumann})
\item Lecture : L. C. Evans, \textit{Partial Differential Equations}, chap. 6
\item Lecture : H. Brezis, \textit{Functional Analysis, Sobolev Spaces and Partial Differential Equations}, chap. 6
\end{refs}

\bigskip

\begin{problem}[{Cauchy--Kowalevski, et l'équation de Lewy sans solution --- Master, short}]
\textbf{PDE-37.} \textbf{Problème.} Énoncer le théorème de Cauchy--Kowalevski --- un problème de
Cauchy analytique à données analytiques sur une surface non caractéristique admet une unique
solution analytique au voisinage de cette surface --- puis montrer qu'il ne peut être renforcé en un
énoncé de bonne position : pour $u_{xx} + u_{yy} = 0$ avec $u(x,0) = 0$ et
$u_y(x,0) = k^{-1}\sin(kx)$, vérifier que $u(x,y) = k^{-2}\sin(kx)\sinh(ky)$, et en conclure
que des données tendant uniformément vers zéro quand $k \to \infty$ produisent des solutions qui
explosent pour tout $y > 0$ fixé. Étudier ensuite l'opérateur de Hans Lewy
\[ Lu = u_{x_1} + i\,u_{x_2} - 2i\,(x_1 + i x_2)\,u_{x_3}, \]
énoncer précisément ce qu'il a démontré en 1957 --- il existe une fonction $f$ lisse pour laquelle
$Lu = f$ n'a de solution dans aucun ouvert --- et expliquer pourquoi cela fait de l'existence dans
la catégorie lisse une question complètement différente de l'existence dans la catégorie
analytique.
\end{problem}

\noindent Sofia Kovalevskaïa a démontré le théorème dans la généralité qui porte son nom
dans sa thèse de 1875, écrite sous la direction de Weierstrass et affinant un résultat de Cauchy de
1842 ; elle fut la première femme d'Europe moderne à recevoir un doctorat en mathématiques, et la
même thèse contient l'exemple montrant que le problème de Cauchy pour l'équation de la chaleur ne
peut être résolu dans les classes analytiques pour des données analytiques arbitraires. Pendant un
demi-siècle, le théorème a paru clore la question de l'existence. L'exemple de Hadamard --- une
partie de ce problème --- a montré que l'analyticité achète l'existence sans la stabilité, de sorte
que le problème de Cauchy pour une équation elliptique est mal posé et qu'aucun schéma numérique ne
peut le sauver ; et la note de trois pages de Lewy dans les \textit{Annals}, qui fut un véritable
choc, a montré que dans la catégorie lisse un opérateur d'apparence parfaitement innocente à
coefficients polynomiaux peut n'avoir aucune solution. De là est née la théorie de la résolubilité
locale, des conditions nécessaires de Hörmander à la condition de Nirenberg--Treves.

\begin{refs}
\item Contexte : Théorème de Cauchy-Kowalevski --- Wikipédia (\url{https://fr.wikipedia.org/wiki/Th%C3%A9or%C3%A8me_de_Cauchy-Kowalevski})
\item Contexte : Exemple de Lewy --- Wikipédia (\url{https://fr.wikipedia.org/wiki/Exemple_de_Lewy})
\item Article : H. Lewy, \og{} An example of a smooth linear partial differential equation without solution \fg{}, \textit{Annals of Mathematics} 66 (1957), 155--158
\item Lecture : L. C. Evans, \textit{Partial Differential Equations}, chap. 4
\end{refs}

\bigskip


\section*{Recherche}

\begin{problem}[{La régularité de Navier--Stokes --- Recherche}]
\textbf{PDE-18.} \textbf{Problème.} Les équations de Navier--Stokes incompressibles
régissent un fluide visqueux :
\[ u_t + (u \cdot \nabla)u = \nu\,\Delta u - \nabla p, \qquad \nabla \cdot u = 0. \]
La question du prix du millénaire : étant données des conditions initiales lisses à décroissance
rapide dans $\mathbb{R}^3$, une solution lisse existe-t-elle pour tout temps --- ou peut-elle
exploser ? \textbf{Statut : ouvert.} Une porte d'entrée vers le pourquoi :
\begin{enumerate}
\item[(a)] démontrer l'identité d'énergie pour les solutions lisses
décroissantes,
\[ \frac{d}{dt}\,\frac{1}{2}\int |u|^2\,dx = -\,\nu \int |\nabla u|^2\,dx, \]
de sorte que l'énergie cinétique totale ne peut que décroître.
\item[(b)] Montrer que les équations ont la
symétrie d'échelle $u_\lambda(x,t) = \lambda\,u(\lambda x, \lambda^2 t)$, et
calculer comment l'énergie de $u_\lambda$ se transforme en dimensions $2$ et $3$.
En conclure qu'en 2-D l'énergie conservée est exactement invariante d'échelle (critique)
tandis qu'en 3-D elle s'affaiblit quand on zoome (surcritique) --- les bornes a priori connues
vivent à la mauvaise échelle pour empêcher une explosion aux petites échelles. Cet écart d'échelle
est l'explication communément admise de la difficulté du problème.
\end{enumerate}

\end{problem}

\noindent Leray a démontré en 1934 que des solutions \textit{faibles} globales existent en
3-D (le concept de PDE-14, inventé exactement à cette fin), mais leur régularité et leur unicité
restent inconnues ; en 2-D, la régularité globale est un théorème (Leray, Ladyjenskaïa). Le
meilleur résultat partiel en 3-D, Caffarelli--Kohn--Nirenberg (1982), montre que l'ensemble des
points singuliers possibles est très petit. Le traitement approfondi de ce problème du prix du
millénaire se trouve sur la page Problèmes ouverts (\url{open-problems.html}).

\begin{refs}
\item Contexte : Solutions des équations de Navier-Stokes --- Wikipédia (\url{https://fr.wikipedia.org/wiki/Solutions_des_%C3%A9quations_de_Navier-Stokes})
\item Article : C. L. Fefferman, \og{} Existence and smoothness of the Navier--Stokes equation \fg{} (2000), description du problème du prix du millénaire, Clay Mathematics Institute
\item Article : J. Leray, \og{} Sur le mouvement d'un liquide visqueux emplissant l'espace \fg{} (1934), Acta Mathematica 63
\end{refs}

\bigskip

\begin{problem}[{L'explosion pour les équations d'Euler incompressibles --- Recherche}]
\textbf{PDE-19.} \textbf{Problème.} Supprimons la viscosité : les équations d'Euler
incompressibles sont $u_t + (u \cdot \nabla)u = -\nabla p$, $\nabla \cdot u = 0$. Des
données initiales lisses et localisées sur $\mathbb{R}^3$ restent-elles lisses pour toujours ?
\textbf{Statut : ouvert sur $\mathbb{R}^3$ pour des données lisses --- mais récemment résolu dans des
cadres voisins, dans le sens de l'explosion.} Elgindi (Annals of Mathematics,
2021) a démontré la formation de singularités en temps fini pour des vitesses de classe
$C^{1,\alpha}$ (une fois dérivables, à dérivée höldérienne --- juste en dessous de la
régularité lisse). Chen et Hou ont ensuite donné une démonstration assistée par ordinateur d'une
explosion en temps fini à partir de données \textit{lisses} dans un domaine cylindrique avec bord,
confirmant rigoureusement le scénario numérique de Luo--Hou (2014) d'un effondrement auto-similaire
au bord ; l'analyse paraît dans un travail en deux parties (2022--2025) et dans les Proceedings
of the National Academy of Sciences (2025). Porte d'entrée :
\begin{enumerate}
\item[(a)] établir l'équation de la
vorticité
\[ \omega_t + (u \cdot \nabla)\,\omega = (\omega \cdot \nabla)\,u, \]
en identifiant $(\omega \cdot \nabla)u$ comme le terme d'étirement des tourbillons.
\item[(b)] Montrer qu'en
2-D le terme d'étirement s'annule identiquement, de sorte que la vorticité est simplement
transportée et que son maximum est conservé --- le mécanisme derrière la régularité globale en
2-D.
\item[(c)] Énoncer le critère de Beale--Kato--Majda (1984) : l'explosion se produit au temps $T$ si
et seulement si $\int_0^T \sup_x |\omega(x,t)|\,dt = \infty$. Expliquer pourquoi tous les
scénarios d'explosion connus sont par conséquent des chasses à la croissance non bornée de la
vorticité.
\end{enumerate}

\end{problem}

\noindent Savoir si un fluide parfait peut spontanément former une singularité est
l'une des plus anciennes questions de la physique mathématique --- les équations sont celles
d'Euler, de 1757. Les années 2020 l'ont fait passer d'une chasse purement numérique à des
théorèmes : la singularité d'Elgindi est entièrement analytique mais exige des données peu
régulières ; celle de Chen--Hou exige un bord mais accepte des données lisses, et sa vérification
rigoureuse par arithmétique d'intervalles est un jalon pour la démonstration assistée par
ordinateur en EDP. Le cas phare --- données lisses, sans bord, tout $\mathbb{R}^3$ --- reste ouvert,
et il est intimement lié à Navier--Stokes (PDE-18) et à l'apparition de la turbulence.

\begin{refs}
\item Contexte : Équations d'Euler (mécanique des fluides) --- Wikipédia (\url{https://fr.wikipedia.org/wiki/%C3%89quations_d%27Euler})
\item Article : T. M. Elgindi, \og{} Finite-time singularity formation for $C^{1,\alpha}$ solutions to the incompressible Euler equations on $\mathbb{R}^3$ \fg{} (2021), Annals of Mathematics 194, arXiv:1904.04795 (\url{https://arxiv.org/abs/1904.04795})
\item Article : J. Chen, T. Y. Hou, \og{} Stable nearly self-similar blowup of the 2D Boussinesq and 3D Euler equations with smooth data \fg{} (2022--2025), arXiv:2210.07191 (\url{https://arxiv.org/abs/2210.07191}) ; voir aussi leur \og{} Singularity formation in 3D Euler equations with smooth initial data and boundary \fg{}, PNAS 122 (2025)
\item Article : J. T. Beale, T. Kato, A. Majda, \og{} Remarks on the breakdown of smooth solutions for the 3-D Euler equations \fg{} (1984), Communications in Mathematical Physics 94
\end{refs}

\bigskip

\begin{problem}[{Rendre rigoureuse la turbulence d'ondes --- Recherche}]
\textbf{PDE-20.} \textbf{Problème.} La théorie de la turbulence d'ondes prédit qu'une mer
d'ondes dispersives faiblement en interaction --- houle de surface, ondes de plasma, ondes de
l'équation de Schrödinger non linéaire --- a un spectre d'énergie moyen $n(k,t)$
évoluant selon une \textit{équation cinétique des ondes}, analogue de l'équation de Boltzmann.
Pour l'équation de Schrödinger non linéaire cubique de dispersion $\omega(k) = |k|^2$, elle
s'écrit
\[ \partial_t n = \iiint \delta(k + k_1 - k_2 - k_3)\,
\delta(\omega + \omega_1 - \omega_2 - \omega_3)\;
n\,n_1 n_2 n_3 \Big( \frac{1}{n} + \frac{1}{n_1} - \frac{1}{n_2} - \frac{1}{n_3} \Big)
\,dk_1\,dk_2\,dk_3, \]
avec $n_i = n(k_i)$, $\omega_i = \omega(k_i)$ : des quadruplets de nombres d'onde
échangent de l'énergie lorsqu'ils satisfont les conditions de résonance
$k + k_1 = k_2 + k_3$ et $\omega + \omega_1 = \omega_2 + \omega_3$. La
dérivation physique (Peierls 1929 ; Hasselmann 1962 ; Zakharov dans les années 1960) est
heuristique.
\textbf{Statut : rigoureusement justifiée dans un cas modèle central ; ouverte en général.}
Deng et Hani (Inventiones Mathematicae, 2023) ont démontré la dérivation complète de cette
équation à partir de l'équation de Schrödinger non linéaire cubique à données initiales
aléatoires, jusqu'à l'échelle de temps cinétique --- l'analogue ondulatoire du théorème de Lanford
pour Boltzmann. Restent ouverts : les cas physiquement premiers (ondes de gravité et de
capillarité à la surface de l'eau, dont l'équation cinétique sous-tend la prévision quotidienne
des états de mer), le comportement au-delà du temps cinétique, et le statut rigoureux des spectres
en cascade de Kolmogorov--Zakharov. Porte d'entrée :
\begin{enumerate}
\item[(a)] pour $\omega(k) = |k|^2$, décrire géométriquement les quadruplets
résonnants.
\item[(b)] Démontrer que les équilibres de Rayleigh--Jeans
$n(k) = T/(\omega(k) + \mu)$ annulent ponctuellement la fonction sous l'intégrale de collision
sur l'ensemble de résonance.
\end{enumerate}

\end{problem}

\noindent L'intuition de Zakharov (honorée par la médaille Dirac) fut que les
équations cinétiques pour les ondes admettent des états stationnaires en loi de puissance --- des
cascades turbulentes --- à côté des équilibres de (b). Faire de tout cela des mathématiques fut un
rêve pendant cinquante ans ; la démonstration de Deng--Hani, appuyée sur des outils venus des EDP
dispersives, des probabilités et même de la combinatoire des diagrammes de Feynman, a ouvert le
domaine. Étendre la rigueur aux ondes de surface, où l'équation de Hasselmann est exploitée
opérationnellement chaque jour par les agences météorologiques, est l'un des défis majeurs les
plus nets à la frontière entre analyse et physique.

\begin{refs}
\item Contexte : Turbulence d'ondes --- Wikipédia (en anglais) (\url{https://en.wikipedia.org/wiki/Wave_turbulence})
\item Article : Y. Deng, Z. Hani, \og{} Full derivation of the wave kinetic equation \fg{} (2023), Inventiones Mathematicae 233, arXiv:2104.11204 (\url{https://arxiv.org/abs/2104.11204})
\item Lecture : S. Nazarenko, \textit{Wave Turbulence}, Springer Lecture Notes in Physics
\end{refs}

\bigskip

\begin{problem}[{L'exposant un tiers d'Onsager --- Recherche, short}]
\textbf{PDE-28.} \textbf{Problème.} Onsager a conjecturé en 1949 qu'une solution
faible des équations d'Euler incompressibles conserve l'énergie cinétique si elle est
höldérienne d'exposant $\alpha > 1/3$, et qu'il existe des solutions faibles dissipant
l'énergie pour tout $\alpha < 1/3$. \textbf{Statut : réglé} --- la moitié
\og{} conservation \fg{} a été démontrée par Constantin, E et Titi (1994) et la moitié
\og{} construction \fg{} par Isett (2018) ; énoncer les deux moitiés précisément, démontrer la moitié
\og{} conservation \fg{}, et expliquer pourquoi $1/3$ est exactement l'exposant que prédit le
scaling de Kolmogorov de 1941, $\delta u(\ell) \sim (\varepsilon\ell)^{1/3}$.
\end{problem}

\noindent Onsager a mis cela dans un unique paragraphe d'un article de 1949 sur
l'hydrodynamique statistique, et cela est resté à peu près non lu pendant quarante ans. L'énoncé
dit quelque chose dont les physiciens avaient besoin : les fluides turbulents dissipent de
l'énergie à un taux qui ne s'annule pas quand la viscosité tend vers zéro, si bien que
l'écoulement non visqueux limite doit être assez irrégulier pour dissiper tout seul. La moitié
\og{} construction \fg{} a fini par être atteinte par intégration convexe, une technique descendant du
théorème de plongement isométrique $C^1$ de Nash et adaptée aux fluides par De Lellis et
Székelyhidi --- la machinerie même qui produit les résultats de non-unicité qui hantent PDE-18.
Noter ce qui n'est \textit{pas} réglé : savoir si les solutions de Navier--Stokes présentent
effectivement une dissipation anomale dans la limite de viscosité évanescente, ce qui est
l'énoncé physique visé par Onsager.

\begin{refs}
\item Contexte : Équations d'Euler (mécanique des fluides) --- Wikipédia (\url{https://fr.wikipedia.org/wiki/%C3%89quations_d%27Euler})
\item Contexte : Turbulence --- Wikipédia (\url{https://fr.wikipedia.org/wiki/Turbulence})
\item Article : P. Isett, \og{} A proof of Onsager's conjecture \fg{} (2018), Annals of Mathematics 188, arXiv:1608.08301 (\url{https://arxiv.org/abs/1608.08301})
\item Article : P. Constantin, W. E, E. S. Titi, \og{} Onsager's conjecture on the energy conservation for solutions of Euler's equation \fg{} (1994), Communications in Mathematical Physics 165
\end{refs}

\bigskip

\begin{problem}[{Où se trouve le point le plus chaud ? --- Recherche, short}]
\textbf{PDE-29.} \textbf{Problème.} La conjecture des points chauds de Rauch (1974) affirme
que pour un domaine borné, la seconde fonction propre de Neumann du laplacien --- c'est-à-dire
le profil de température en temps long d'un corps parfaitement isolé --- atteint son
maximum et son minimum uniquement sur la frontière, de sorte que le point le plus chaud finit par
migrer vers le bord. \textbf{Statut : fausse en général, ouverte dans les cas qui comptent} ---
des contre-exemples sont connus pour des domaines plans à trous (Burdzy--Werner 1999 ;
Burdzy 2005) et pour des domaines convexes en dimension suffisamment grande
(de Dios Pont, 2024), tandis que la conjecture est un théorème pour tout triangle
euclidien (Judge--Mondal 2020, avec un rectificatif de 2022) ; elle reste ouverte pour les domaines
plans simplement connexes et pour les domaines plans convexes, et votre tâche est d'énoncer
la conjecture assez précisément pour que toutes ces affirmations soient véritablement
distinctes.
\end{problem}

\noindent La conjecture a une jumelle probabiliste : la fonction propre gouverne
l'endroit où un mouvement brownien réfléchi a des chances de se trouver aux grands temps,
ce qui explique pourquoi ce sont les probabilistes qui ont produit à la fois les contre-exemples et
l'essentiel de la théorie positive. C'est un bon exemple d'un énoncé que tout le monde juge évident
et qui reste néanmoins largement non démontré --- un projet Polymath a attaqué le cas du triangle
acutangle en 2012 sans le trancher, et le théorème complet sur les triangles a demandé huit années
de plus et un argument long et délicat.

\begin{refs}
\item Contexte : Condition aux limites de Neumann --- Wikipédia (\url{https://fr.wikipedia.org/wiki/Condition_aux_limites_de_Neumann})
\item Contexte : Géométrie spectrale --- Wikipédia (\url{https://fr.wikipedia.org/wiki/G%C3%A9om%C3%A9trie_spectrale})
\item Article : C. Judge, S. Mondal, \og{} Euclidean triangles have no hot spots \fg{} (2020), Annals of Mathematics 191 ; rectificatif, Annals of Mathematics 195 (2022), arXiv:1802.01800 (\url{https://arxiv.org/abs/1802.01800})
\item Article : K. Burdzy, W. Werner, \og{} A counterexample to the `hot spots' conjecture \fg{} (1999), Annals of Mathematics 149 ; J. de Dios Pont, \og{} Convex sets can have interior hot spots \fg{} (2024), arXiv:2412.06344 (\url{https://arxiv.org/abs/2412.06344})
\end{refs}

\bigskip

\begin{problem}[{L'écoulement quasi-géostrophique de surface surcritique --- Recherche}]
\textbf{PDE-30.} \textbf{Problème.} L'équation quasi-géostrophique de surface (SQG) pour un scalaire
$\theta(x,t)$ sur $\mathbb{R}^2$ est
\[ \theta_t + u\cdot\nabla\theta + \kappa(-\Delta)^{\gamma}\theta = 0, \qquad
u = \nabla^{\perp}(-\Delta)^{-1/2}\theta, \]
de sorte que la vitesse est une transformée intégrale singulière d'ordre zéro de la quantité
transportée --- exactement la relation structurelle qu'entretient la vorticité avec la vitesse dans
l'écoulement d'Euler tridimensionnel, mais avec une dimension de moins.
\textbf{Statut :} pour $\gamma > 1/2$ (sous-critique) la régularité globale est classique ;
pour $\gamma = 1/2$ (critique) elle a été démontrée en 2007, indépendamment et par des méthodes
complètement différentes, par Kiselev--Nazarov--Volberg et par
Caffarelli--Vasseur ; pour $0 < \gamma < 1/2$ (surcritique) et pour le cas non visqueux
$\kappa = 0$, \textbf{la régularité globale pour des données lisses est ouverte}. Y travailler
comme suit.
\begin{enumerate}
\item[(a)] Montrer que si $\theta$ résout l'équation, alors
$\theta_\lambda(x,t) = \lambda^{2\gamma-1}\theta(\lambda x, \lambda^{2\gamma}t)$ aussi,
et en déduire que $\gamma = 1/2$ est précisément l'exposant pour lequel le changement d'échelle
laisse la norme $L^\infty$ invariante. Expliquer pourquoi cela rend $\gamma = 1/2$
critique, étant donné que $L^\infty$ est la norme la plus forte que l'équation soit connue pour
contrôler.
\item[(b)] Démontrer le principe du maximum $L^p$ : pour $\kappa \ge 0$ et tout
$p \in [1,\infty]$, $\|\theta(\cdot,t)\|_{L^p}$ est décroissante en $t$.
\item[(c)] Démontrer le critère d'explosion de Constantin--Majda--Tabak pour SQG non visqueuse : une
solution lisse persiste sur $[0,T]$ pourvu que
$\int_0^T \|\nabla^{\perp}\theta(\cdot,t)\|_{L^\infty}\,dt < \infty$. Le comparer
ligne à ligne avec le critère de Beale--Kato--Majda pour l'écoulement d'Euler tridimensionnel
et dire ce qui y joue le rôle de la vorticité.
\item[(d)] Étudier les deux démonstrations du cas critique --- l'argument non local de module de
continuité et l'itération de De Giorgi --- et identifier exactement
quelle inégalité échoue dans chacune lorsque $\gamma < 1/2$.
\end{enumerate}

\end{problem}

\noindent Constantin, Majda et Tabak ont proposé SQG en 1994 comme
laboratoire bidimensionnel du problème d'explosion d'Euler en dimension trois de
PDE-19 : l'analogie avec l'étirement de la vorticité n'est pas vague mais formelle, terme à terme,
et une singularité ici serait une indication forte sur ce qui se passe là-bas. C'est aussi de la
vraie physique --- l'équation régit les fronts de température à la surface d'un fluide stratifié en
rotation rapide, c'est-à-dire les fronts météorologiques d'une carte synoptique.
La double résolution du cas critique en 2007 est l'un des meilleurs cas d'étude, dans l'analyse
moderne, de deux routes indépendantes vers un même théorème, et la plage surcritique est là où se
mène le combat actuel.

\begin{refs}
\item Contexte : Approximation quasi-géostrophique --- Wikipédia (\url{https://fr.wikipedia.org/wiki/Approximation_quasi-g%C3%A9ostrophique})
\item Article : P. Constantin, A. J. Majda, E. Tabak, \og{} Formation of strong fronts in the 2-D quasigeostrophic thermal active scalar \fg{} (1994), Nonlinearity 7
\item Article : A. Kiselev, F. Nazarov, A. Volberg, \og{} Global well-posedness for the critical 2D dissipative quasi-geostrophic equation \fg{} (2007), Inventiones Mathematicae 167, arXiv:math/0604185 (\url{https://arxiv.org/abs/math/0604185})
\item Article : L. Caffarelli, A. Vasseur, \og{} Drift diffusion equations with fractional diffusion and the quasi-geostrophic equation \fg{} (2010), Annals of Mathematics 171, arXiv:math/0608447 (\url{https://arxiv.org/abs/math/0608447})
\end{refs}

\bigskip

\begin{problem}[{La conjecture de De Giorgi pour l'équation d'Allen--Cahn --- Recherche}]
\textbf{PDE-38.} \textbf{Problème.} On considère les solutions bornées de l'équation d'Allen--Cahn
\[ \Delta u + u - u^3 = 0 \quad \text{sur } \mathbb{R}^n, \qquad |u| \le 1, \qquad
\frac{\partial u}{\partial x_n} > 0 . \]
De Giorgi a conjecturé en 1978 que toute solution de ce type est unidimensionnelle --- c'est-à-dire
$u(x) = g(x \cdot e + b)$ pour un certain vecteur unitaire $e$ --- au moins lorsque
$n \le 8$.
\begin{enumerate}
\item[(a)] Traiter complètement le cas unidimensionnel. Démontrer que les solutions bornées de
$u'' = u^3 - u$ sur $\mathbb{R}$ avec $u(\pm\infty) = \pm 1$ sont exactement les translatées
de $\tanh(x/\sqrt{2})$, en établissant l'intégrale première
$\tfrac{1}{2}(u')^2 = \tfrac{1}{4}(1-u^2)^2$ et en l'intégrant. Vérifier ensuite que
$u(x) = \tanh\bigl((x\cdot e + b)/\sqrt{2}\bigr)$ résout l'équation en toute
dimension.
\item[(b)] Identifier la structure variationnelle : montrer que l'équation est l'équation
d'Euler--Lagrange de
\[ E(u) = \int \Bigl( \tfrac{1}{2}|\nabla u|^2 + \tfrac{1}{4}(1-u^2)^2 \Bigr) dx , \]
et décrire le théorème de Modica et Mortola : après changement d'échelle, ces énergies convergent
(au sens de la $\Gamma$-convergence) vers une constante fois l'aire de l'interface entre les
phases $u = 1$ et $u = -1$. Expliquer pourquoi cela fait que les lignes de niveau d'une
solution se comportent comme une surface minimale, et pourquoi la conjecture est donc un énoncé
déguisé sur la rigidité des surfaces minimales.
\item[(c)] Rendre l'analogie exacte. Énoncer le problème de Bernstein --- un graphe minimal entier sur
$\mathbb{R}^m$ est un hyperplan exactement lorsque $m \le 7$, le contre-exemple pour
$m \ge 8$ étant construit par Bombieri, De Giorgi et Giusti (1969) à partir du cône de Simons ---
et expliquer pourquoi le seuil $m \le 7$ là-bas correspond à $n \le 8$ ici.
\item[(d)] Rendre compte du statut, précisément. La conjecture est un théorème pour $n = 2$
(Ghoussoub--Gui, 1998) et $n = 3$ (Ambrosio--Cabré, 2000) ; Savin l'a démontrée pour
$4 \le n \le 8$ sous l'hypothèse supplémentaire que $u(x', x_n) \to \pm 1$ ponctuellement quand
$x_n \to \pm\infty$ (2009) ; et del Pino, Kowalczyk et Wei ont construit un contre-exemple pour
tout $n \ge 9$ (2011) en recollant le profil de la partie (a) le long du graphe de
Bombieri--De Giorgi--Giusti. Expliquer ce que l'hypothèse supplémentaire du théorème de Savin exclut,
et pourquoi la supprimer est difficile.
\textit{Statut : ouvert pour $4 \le n \le 8$ sans l'hypothèse de limite ; réglé sinon --- vrai pour
$n \le 3$, faux pour $n \ge 9$.}
\end{enumerate}

\end{problem}

\noindent Ennio De Giorgi a proposé cela en 1978, et il était la bonne personne pour le
faire : il avait démontré en 1957 la régularité höldérienne qui a réglé le dix-neuvième problème de
Hilbert, et il avait construit, avec Bombieri et Giusti, le contre-exemple de graphe minimal
même qui marque aujourd'hui la limite de sa propre conjecture. Que les deux seuils dimensionnels
s'alignent --- le cône de Simons apparaît en dimension huit, la transition de phase ploie en
dimension neuf --- est l'une des correspondances les plus frappantes de l'analyse, et c'est la raison
pour laquelle la conjecture est vue comme le miroir, côté transition de phase, du problème de
Bernstein. Le contre-exemple de del Pino--Kowalczyk--Wei mérite d'être étudié comme construction en
soi : prendre un objet géométrique, recoller le profil unidimensionnel le long de celui-ci, puis
corriger l'erreur en une solution exacte.

\begin{refs}
\item Contexte : Équation d'Allen-Cahn --- Wikipédia (\url{https://fr.wikipedia.org/wiki/%C3%89quation_d%27Allen-Cahn})
\item Contexte : Problème de Bernstein --- Wikipédia (\url{https://fr.wikipedia.org/wiki/Probl%C3%A8me_de_Bernstein})
\item Article : O. Savin, \og{} Regularity of flat level sets in phase transitions \fg{}, \textit{Annals of Mathematics} 169 (2009), 41--78
\item Article : M. del Pino, M. Kowalczyk \& J. Wei, \og{} On De Giorgi's conjecture in dimension $N \ge 9$ \fg{}, \textit{Annals of Mathematics} 174 (2011), 1485--1569, arXiv:0806.3141 (\url{https://arxiv.org/abs/0806.3141})
\end{refs}

\bigskip

\begin{problem}[{Quelle est la longueur d'un ensemble nodal ? La conjecture de Yau --- Recherche}]
\textbf{PDE-39.} \textbf{Problème.} Soit $(M,g)$ une variété riemannienne lisse fermée de dimension $n$ et
soit $\varphi_\lambda$ vérifiant $\Delta \varphi_\lambda + \lambda \varphi_\lambda = 0$. Son
\textit{ensemble nodal} est $Z = \{\varphi_\lambda = 0\}$, le motif d'onde stationnaire que
Chladni a rendu visible avec du sable. Yau a conjecturé que
\[ c\,\sqrt{\lambda} \;\le\; \mathcal{H}^{n-1}(Z) \;\le\; C\,\sqrt{\lambda}, \]
avec des constantes ne dépendant que de $(M,g)$.
\begin{enumerate}
\item[(a)] Vérifier la conjecture là où le calcul est possible. Sur le tore plat
$\mathbb{R}^2/(2\pi\mathbb{Z})^2$, prendre $\varphi = \sin(mx)\sin(ny)$, avec
$\lambda = m^2 + n^2$ ; décrire l'ensemble nodal, calculer exactement sa longueur totale, et
vérifier qu'elle est encadrée par des multiples constants de $\sqrt{\lambda}$. Faire le même
décompte qualitatif pour une harmonique sphérique de degré $\ell$ sur la sphère ronde, où
$\lambda = \ell(\ell+1)$.
\item[(b)] Démontrer la moitié élémentaire de l'intuition. Montrer qu'une solution de l'équation de
Helmholtz oscille à l'échelle de longueur $\lambda^{-1/2}$, énoncer le théorème des domaines
nodaux de Courant (la $k$-ième fonction propre a au plus $k$ domaines nodaux) et le démontrer
dans le cas unidimensionnel de Sturm--Liouville, et le combiner avec l'inégalité de Faber--Krahn
pour montrer que tout domaine nodal a un volume au moins $c\lambda^{-n/2}$. Expliquer pourquoi
cela rend plausible une minoration de l'ensemble nodal sans pour autant la démontrer.
\item[(c)] Étudier le cas analytique. Énoncer le théorème de Donnelly et Fefferman (1988) : les deux
bornes valent lorsque la métrique $g$ est réelle-analytique. Décrire l'instrument qui rend ce
cas abordable --- l'indice de doublement, ou fonction de fréquence, d'Almgren puis de Garofalo et
Lin --- et expliquer ce qu'il mesure.
\item[(d)] Rendre compte du statut moderne. Alexander Logunov a démontré la minoration
$c\sqrt{\lambda}$ pour les métriques lisses, et dans un article compagnon une majoration
$\mathcal{H}^{n-1}(Z) \le C\lambda^{\alpha}$ pour un certain exposant $\alpha(n) > 1/2$ ;
expliquer pourquoi la majoration est la moitié la plus difficile, et ce qu'une démonstration de
l'exposant optimal devrait contrôler.
\textit{Statut : la minoration de la conjecture de Yau est un théorème pour les métriques lisses
(Logunov, Annals of Mathematics, 2018) ; la majoration optimale $C\sqrt{\lambda}$ est ouverte en
toute dimension $n \ge 2$ pour les métriques lisses --- en dimension deux, les meilleures bornes
connues sont de la forme $C\lambda^{3/4}$, avec de légères améliorations depuis --- tandis que
pour les métriques réelles-analytiques les deux bornes ont été réglées en 1988.}
\end{enumerate}

\end{problem}

\noindent Yau a inclus cette question dans sa liste de problèmes ouverts de géométrie de
1982, et c'est la question quantitative la plus grossière que l'on puisse poser sur les images de
cette page : quelle quantité d'ensemble nodal y a-t-il ? La réponse est forcément d'ordre
$\sqrt{\lambda}$ parce qu'une fonction propre oscille à l'échelle $\lambda^{-1/2}$, et pourtant
transformer cette heuristique en théorème pour des métriques lisses générales a demandé
trente-cinq ans. La démonstration de Logunov, mise au point avec Eugenia Malinnikova, propage
l'indice de doublement à travers les échelles par un argument combinatoire élémentaire dans ses
ingrédients et redoutable dans sa comptabilité ; elle a depuis été adaptée à plusieurs autres
problèmes d'analyse des fonctions propres. Le sujet se relie directement au reste de cette page ---
les motifs nodaux de PDE-24 et PDE-31, la géométrie spectrale de PDE-17, et les fonctions propres
de Neumann de PDE-29.

\begin{refs}
\item Contexte : Géométrie spectrale --- Wikipédia (\url{https://fr.wikipedia.org/wiki/G%C3%A9om%C3%A9trie_spectrale})
\item Article : H. Donnelly \& C. Fefferman, \og{} Nodal sets of eigenfunctions on Riemannian manifolds \fg{}, \textit{Inventiones Mathematicae} 93 (1988), 161--183
\item Article : A. Logunov, \og{} Nodal sets of Laplace eigenfunctions: proof of Nadirashvili's conjecture and of the lower bound in Yau's conjecture \fg{}, \textit{Annals of Mathematics} 187 (2018), 241--262, arXiv:1605.02589 (\url{https://arxiv.org/abs/1605.02589})
\item Article : A. Logunov, \og{} Nodal sets of Laplace eigenfunctions: polynomial upper estimates of the Hausdorff measure \fg{}, \textit{Annals of Mathematics} 187 (2018), 221--239, arXiv:1605.02587 (\url{https://arxiv.org/abs/1605.02587})
\end{refs}

\bigskip

\begin{problem}[{La conjecture de Schiffer : quel tambour est un disque ? --- Recherche, short}]
\textbf{PDE-40.} \textbf{Problème.} La conjecture de Schiffer affirme que si un domaine borné
$\Omega \subset \mathbb{R}^n$ à frontière lipschitzienne connexe admet une solution non
constante de $\Delta u + \lambda u = 0$ avec $\lambda > 0$ satisfaisant \textit{à la fois}
$\partial u/\partial\nu = 0$ et $u = \text{const}$ sur $\partial\Omega$, alors $\Omega$ est
une boule. Démontrer que le disque admet bien une telle fonction --- prendre $u(r) = J_0(kr)$ avec
$ka$ un zéro de $J_1$, et utiliser le fait que $J_0$ et $J_1$ n'ont aucun zéro positif
commun (hypothèse de Bourget, démontrée par Siegel en 1929) pour voir que $u$ est une constante
non nulle sur le bord --- puis énoncer la formulation équivalente en termes du problème de Pompeiu
et rendre compte honnêtement de ce qui est connu et de ce qui ne l'est pas.
\end{problem}

\noindent Dimitrie Pompeiu a demandé en 1929 si un domaine plan dont les intégrales sur
tous ses déplacements annulent une fonction continue force cette fonction à être nulle ; les
domaines où cela échoue sont les domaines exceptionnels, et le disque est exceptionnel. S. A.
Williams a montré en 1976 qu'un domaine exceptionnel doit porter exactement la fonction propre de
Neumann surdéterminée ci-dessus, ce qui relie la question de Pompeiu en géométrie intégrale à une
question d'EDP attribuée à Menahem Schiffer. Le contraste instructif est avec le théorème de Serrin
de 1971 : pour $\Delta u = -1$ avec $u = 0$ et $\partial u/\partial \nu$ constante sur la
frontière, le domaine \textit{doit} être une boule, ce qui se démontre par la méthode du plan mobile
d'Alexandrov --- mais cet argument a besoin d'un signe, et une équation aux valeurs propres n'en a
pas. \textit{Statut : ouvert. Des résultats partiels sont connus sous des hypothèses supplémentaires de
régularité, de symétrie ou de convexité, et le problème compte parmi les problèmes ouverts
surdéterminés les plus connus des EDP elliptiques.}

\begin{refs}
\item Contexte : Problème de Pompeiu (et conjecture de Schiffer) --- Wikipédia (\url{https://fr.wikipedia.org/wiki/Probl%C3%A8me_de_Pompeiu})
\item Article : S. A. Williams, \og{} A partial solution of the Pompeiu problem \fg{}, \textit{Mathematische Annalen} 223 (1976), 183--190
\item Article : J. Serrin, \og{} A symmetry problem in potential theory \fg{}, \textit{Archive for Rational Mechanics and Analysis} 43 (1971), 304--318
\item Lecture : L. E. Fraenkel, \textit{An Introduction to Maximum Principles and Symmetry in Elliptic Problems}
\end{refs}

\bigskip


\section*{Pour aller plus loin}


\vfill
\noindent\emph{Hors de la Caverne} --- des probl\`emes, pas de r\'eponses.
\end{document}
